\documentclass[12pt]{amsart}

\usepackage{amsmath,amssymb,amsthm}
\usepackage{mathrsfs}
\usepackage{mathtools}
\usepackage{tikz-cd}
\usepackage{enumitem}
\usepackage{hyperref}
\usepackage{cleveref}

% Theorem environments (shared counter within section)
\newtheorem{theorem}{Theorem}[section]
\newtheorem{proposition}[theorem]{Proposition}
\newtheorem{lemma}[theorem]{Lemma}
\newtheorem{corollary}[theorem]{Corollary}

\theoremstyle{definition}
\newtheorem{definition}[theorem]{Definition}
\newtheorem{example}[theorem]{Example}

\theoremstyle{remark}
\newtheorem*{remark}{Remark}

% Notation
\DeclareMathOperator{\Var}{Var}
\DeclareMathOperator{\tr}{tr}
\newcommand{\R}{\mathbb{R}}
\newcommand{\Rpos}{\mathbb{R}_{>0}}
\newcommand{\bT}{\mathbb{T}}
\newcommand{\Z}{\mathbb{Z}}
\newcommand{\cN}{\mathcal{N}}
\newcommand{\cF}{\mathcal{F}}
\newcommand{\cM}{\mathcal{M}}
\newcommand{\cP}{\mathcal{P}}
\newcommand{\Htwo}{H^2}
\newcommand{\cK}{\mathcal{K}}
\newcommand{\cT}{\mathcal{T}}
\DeclareMathOperator{\ind}{ind}
\newcommand{\Proj}{\mathsf{Proj}}

\title[Additive and Multiplicative CLT]{Additive and multiplicative central limit theorems:\\ a duality via Fourier analysis on groups}
\author{Jianghan Zhang}
\address{Department of Mathematics, Dartmouth College, Hanover, NH 03755}
\email{jianghan.zhang.gr@dartmouth.edu}
\date{\today}

\subjclass[2020]{60F05, 43A25, 60B15, 30H10, 47B35, 22E30}
\keywords{Central limit theorem, locally compact abelian groups, Mellin transform, Gaussian measures, heat semigroup, Hardy spaces, Toeplitz operators, Wold decomposition, Szeg\H{o} limit theorem, Peter-Weyl theorem}

\begin{document}

\begin{abstract}
We present a unified treatment of the additive and multiplicative central limit theorems as instances of a single phenomenon on locally compact abelian groups. The classical CLT on $(\R, +)$ and the multiplicative CLT on $(\Rpos, \times)$ are shown to be related by the Fourier--Mellin duality, which arises from the group isomorphism $\log : (\Rpos, \times) \to (\R, +)$. We develop this correspondence at the level of characteristic functions, heat equations, and Gaussian measures, and place it in the general framework of harmonic analysis on LCA groups. We indicate how the Peter-Weyl decomposition extends the spectral theory to compact non-abelian groups, where the Casimir eigenvalues govern irreversible convergence to equilibrium.
\end{abstract}

\maketitle

%% ============================================================
\section*{Introduction}

The classical central limit theorem asserts that normalized sums of independent, identically distributed random variables converge in distribution to a Gaussian. Less commonly discussed is its multiplicative counterpart: normalized products of positive i.i.d.\ random variables converge to a log-normal distribution. These two results are usually proved separately, with the multiplicative version dismissed as a trivial consequence of applying the logarithm.

The purpose of this note is to show that this ``trivial'' reduction is, in fact, an instance of a deep structural duality. The additive and multiplicative CLTs live on different groups---$(\R, +)$ and $(\Rpos, \times)$, respectively---and the logarithm $\log : (\Rpos, \times) \to (\R, +)$ is an isomorphism of locally compact abelian (LCA) groups. The Fourier transform on $(\R, +)$ and the Mellin transform on $(\Rpos, \times)$ are related by this isomorphism, yielding a commutative diagram that encodes the duality precisely.

We develop this perspective in detail, showing how characteristic functions, heat equations, and Gaussian measures transform under the Fourier--Mellin correspondence. We then place these results in the general framework of harmonic analysis on LCA groups, where the ``Gaussian'' on each group is characterized as the unique \emph{probability density} that is a fixed point of the Fourier transform and as the fundamental solution of the associated heat equation.

\medskip
\noindent\textbf{Notation.}
Throughout, $\R$ denotes the reals, $\Rpos = (0, \infty)$ the positive reals, $\bT = \R / \Z$ the circle group, and $\Z$ the integers. For a topological group $G$, we write $\widehat{G}$ for its Pontryagin dual. The Fourier transform (characteristic function) of a probability measure $\mu$ on $G$ is denoted $\widehat{\mu}$. We write $\cF$ for the Fourier transform on $(\R, +)$ and $\cM$ for the Mellin transform on $(\Rpos, \times)$.

%% ============================================================
\section{Foundations of probability theory}\label{sec:probability}

We recall the measure-theoretic foundations of probability, following Kolmogorov's axiomatization \cite{Billingsley,Durrett}. The purpose of this section is twofold: to fix notation used throughout, and to exhibit the structures---$\sigma$-algebras, measures, conditional expectations---that the Fourier--Mellin duality will later transform.

\begin{definition}[Probability space]\label{def:prob-space}
A \emph{probability space} is a triple $(\Omega, \cF, P)$ consisting of:
\begin{enumerate}[label=(\roman*)]
\item a non-empty set $\Omega$ (the \emph{sample space}),
\item a $\sigma$-algebra $\cF \subseteq 2^\Omega$ (the \emph{event algebra}),
\item a measure $P : \cF \to [0,1]$ with $P(\Omega) = 1$ (the \emph{probability measure}).
\end{enumerate}
An element $\omega \in \Omega$ is an \emph{outcome}; a set $A \in \cF$ is an \emph{event}.
\end{definition}

\begin{remark}
The axioms encode three intuitive requirements. First, $\cF$ is closed under countable set operations, so any statement built from countably many events remains an event. Second, $P(\varnothing) = 0$ and $P(\Omega) = 1$, normalizing certainty. Third, $P$ is \emph{countably additive}: for pairwise disjoint events $A_1, A_2, \ldots \in \cF$,
\[
P\!\left(\bigcup_{n=1}^\infty A_n\right) = \sum_{n=1}^\infty P(A_n).
\]
Countable additivity is the key axiom; it is what makes the theory non-trivial, distinguishing probability from finitely additive ``charges'' and enabling limiting arguments such as the CLT.
\end{remark}

\begin{definition}[Random variable]\label{def:rv}
A \emph{random variable} on $(\Omega, \cF, P)$ is a measurable function $X : (\Omega, \cF) \to (\R, \mathcal{B}(\R))$, where $\mathcal{B}(\R)$ denotes the Borel $\sigma$-algebra. More generally, for a measurable space $(S, \mathcal{S})$, an \emph{$S$-valued random variable} is a measurable map $X : (\Omega, \cF) \to (S, \mathcal{S})$. The \emph{law} (or \emph{distribution}) of $X$ is the pushforward measure $\mu_X = P \circ X^{-1}$ on $(S, \mathcal{S})$.
\end{definition}

\begin{definition}[Conditional probability]\label{def:cond-prob}
Let $(\Omega, \cF, P)$ be a probability space and let $B \in \cF$ with $P(B) > 0$. The \emph{conditional probability} of $A \in \cF$ given $B$ is
\[
P(A \mid B) \;=\; \frac{P(A \cap B)}{P(B)}.
\]
\end{definition}

\begin{remark}
Conditioning on an event $B$ with $P(B) > 0$ restricts the probability space: the map $A \mapsto P(A \mid B)$ is itself a probability measure on $(\Omega, \cF)$, concentrated on $B$. The definition encodes a selection effect---knowing that $B$ has occurred reweights the sample space.
\end{remark}

\begin{proposition}[Law of total probability]\label{prop:total-prob}
Let $B_1, B_2, \ldots$ be a countable partition of $\Omega$ into events of positive probability. Then for any $A \in \cF$,
\[
P(A) = \sum_{n=1}^\infty P(A \mid B_n)\, P(B_n).
\]
\end{proposition}

\begin{proof}
Since $A = \bigcup_n (A \cap B_n)$ with the sets $A \cap B_n$ pairwise disjoint,
\[
P(A) = \sum_n P(A \cap B_n) = \sum_n P(A \mid B_n)\, P(B_n). \qedhere
\]
\end{proof}

\begin{theorem}[Bayes' rule]\label{thm:bayes}
Let $A, B \in \cF$ with $P(A) > 0$ and $P(B) > 0$. Then
\begin{equation}\label{eq:bayes}
P(A \mid B) \;=\; \frac{P(B \mid A)\, P(A)}{P(B)}.
\end{equation}
More generally, if $\{A_n\}$ is a countable partition of $\Omega$ with $P(A_n) > 0$ for each $n$, then
\begin{equation}\label{eq:bayes-partition}
P(A_k \mid B) \;=\; \frac{P(B \mid A_k)\, P(A_k)}{\displaystyle\sum_n P(B \mid A_n)\, P(A_n)}.
\end{equation}
\end{theorem}

\begin{proof}
Equation~\eqref{eq:bayes} follows from the symmetry $P(A \cap B) = P(B \mid A)\, P(A) = P(A \mid B)\, P(B)$. Substituting the law of total probability (\cref{prop:total-prob}) into the denominator yields \eqref{eq:bayes-partition}.
\end{proof}

\begin{remark}[Bayesian interpretation]
Bayes' rule inverts the direction of conditioning. In the Bayesian reading, $P(A_k)$ is the \emph{prior} probability of hypothesis $A_k$; $P(B \mid A_k)$ is the \emph{likelihood} of the observed data $B$ under hypothesis $A_k$; and $P(A_k \mid B)$ is the \emph{posterior} probability of $A_k$ given the data. The denominator $P(B) = \sum_n P(B \mid A_n)\, P(A_n)$ is a normalizing constant (the \emph{marginal likelihood} or \emph{evidence}). Bayes' rule is thus the mechanism by which data transforms prior beliefs into posterior beliefs.
\end{remark}

\begin{definition}[Expectation]\label{def:expectation}
Let $X$ be a random variable on $(\Omega, \cF, P)$. The \emph{expectation} (or \emph{mean}) of $X$ is the Lebesgue integral
\[
E[X] = \int_\Omega X(\omega)\, dP(\omega),
\]
provided $\int_\Omega |X|\, dP < \infty$ (i.e., $X \in L^1(\Omega, \cF, P)$). Equivalently, by the change-of-variables formula for pushforward measures,
\[
E[X] = \int_\R x\, d\mu_X(x),
\]
where $\mu_X = P \circ X^{-1}$ is the law of $X$.
\end{definition}

\begin{remark}
The expectation is a linear functional on $L^1(P)$: $E[aX + bY] = aE[X] + bE[Y]$. It is also monotone ($X \leq Y$ a.s.\ implies $E[X] \leq E[Y]$) and continuous under dominated convergence. These three properties---linearity, monotonicity, and continuity---characterize the Lebesgue integral and are the workhorses of probabilistic computation.
\end{remark}

\begin{remark}[Variance]
The \emph{variance} $\Var(X) = E\bigl[(X - E[X])^2\bigr] = E[X^2] - (E[X])^2$ measures the spread of $X$ around its mean. It is the quadratic form that appears in the CLT (\cref{thm:add-clt}): the normalization $S_n / \sqrt{n}$ is chosen precisely so that $\Var(S_n / \sqrt{n}) = \sigma^2$ for all $n$.
\end{remark}

We turn to the central construction: conditional expectation with respect to a sub-$\sigma$-algebra. This is the measure-theoretic generalization of the elementary conditional probability and requires a non-trivial existence theorem.

\begin{definition}[Conditional expectation]\label{def:cond-exp}
Let $X \in L^1(\Omega, \cF, P)$ and let $\mathcal{G} \subseteq \cF$ be a sub-$\sigma$-algebra. The \emph{conditional expectation} of $X$ given $\mathcal{G}$, denoted $E[X \mid \mathcal{G}]$, is the $P$-a.s.\ unique $\mathcal{G}$-measurable function satisfying
\begin{equation}\label{eq:cond-exp}
\int_G E[X \mid \mathcal{G}]\, dP = \int_G X\, dP \qquad \text{for all } G \in \mathcal{G}.
\end{equation}
\end{definition}

\begin{theorem}[Existence and uniqueness]\label{thm:cond-exp-exists}
For any $X \in L^1(\Omega, \cF, P)$ and any sub-$\sigma$-algebra $\mathcal{G} \subseteq \cF$, the conditional expectation $E[X \mid \mathcal{G}]$ exists and is unique $P$-a.s.
\end{theorem}

\begin{proof}
Define the signed measure $\nu(G) = \int_G X\, dP$ for $G \in \mathcal{G}$. Then $\nu \ll P|_\mathcal{G}$ (absolute continuity follows from $X \in L^1$). By the Radon--Nikod\'ym theorem, there exists a $\mathcal{G}$-measurable function $f$ with $\nu(G) = \int_G f\, dP$ for all $G \in \mathcal{G}$, unique $P$-a.s. Set $E[X \mid \mathcal{G}] = f$.
\end{proof}

\begin{remark}[Projection interpretation]
The conditional expectation $E[X \mid \mathcal{G}]$ is the orthogonal projection of $X$ onto $L^2(\Omega, \mathcal{G}, P)$ when $X \in L^2$. This geometric interpretation illuminates its fundamental properties:
\begin{enumerate}[label=(\roman*)]
\item \textbf{Tower property}: if $\mathcal{H} \subseteq \mathcal{G} \subseteq \cF$, then $E\bigl[E[X \mid \mathcal{G}] \mid \mathcal{H}\bigr] = E[X \mid \mathcal{H}]$. (Composing two projections onto nested subspaces yields the projection onto the smaller subspace.)
\item \textbf{Pulling out what is known}: if $Y$ is $\mathcal{G}$-measurable and bounded, then $E[XY \mid \mathcal{G}] = Y \cdot E[X \mid \mathcal{G}]$.
\item \textbf{Independence}: if $\sigma(X)$ and $\mathcal{G}$ are independent, then $E[X \mid \mathcal{G}] = E[X]$.
\end{enumerate}
Property~(iii) is the probabilistic content of orthogonality: independent information is uninformative.
\end{remark}

\begin{example}[La Gioconda and the lattice of $\sigma$-algebras]\label{ex:gioconda}
Let $\Omega$ be the set of all paintings entitled ``Mona Lisa'' (La Gioconda)---Leonardo's original at the Louvre, the near-contemporaneous copy at the Museo del Prado, and the countless reproductions produced over five centuries. Equip $\Omega$ with the uniform probability measure and consider three nested sub-$\sigma$-algebras:
\[
\mathcal{H} \;\subseteq\; \mathcal{G} \;\subseteq\; \cF,
\]
where:
\begin{enumerate}[label=(\roman*)]
\item $\mathcal{H} = \sigma(\mathrm{file\; format})$: the coarsest level---an observer sees only that the image is a \texttt{.jpg} file at $960$\,px resolution. At this resolution, Leonardo's original and the Prado copy are \emph{indistinguishable}: they belong to the same $\mathcal{H}$-atom.
\item $\mathcal{G} = \sigma(\mathrm{author})$: the intermediate level---the observer knows who painted each work. Conditioning on the event $\{\mathrm{painter} = \text{Leonardo}\}$ now picks out a single element of $\Omega$. \emph{The definite article ``the'' in ``THE Mona Lisa'' is a conditional operation}: it is well-defined only relative to a $\sigma$-algebra fine enough to separate the original from its copies.
\item $\cF$: full provenance---canvas material, pigment chemistry, exhibition history. Every painting is distinguished.
\end{enumerate}

The tower property reads concretely: if you first identify the painting by authorship ($\mathcal{G}$), then forget the author and retain only the file format ($\mathcal{H}$), you obtain the same answer as if you had consulted only the file format from the start. Information, once discarded, cannot be recovered by re-projection.

The moral is general: a $\sigma$-algebra is not a technicality---it is a \emph{resolution}, a choice of what distinctions the observer can make. Conditional expectation $E[X \mid \mathcal{G}]$ returns the best approximation to $X$ that is visible at the resolution~$\mathcal{G}$. The lattice of sub-$\sigma$-algebras is the lattice of all possible levels of knowledge, and the tower property asserts its internal consistency. In this light, probability theory is not primarily a theory of ``randomness''---it is a theory of \emph{information}. Randomness is not in the phenomenon; it is in the $\sigma$-algebra of the observer.
\end{example}

\begin{example}[Conditional expectation recovers conditional probability]\label{ex:cond-exp-partition}
Let $\{B_n\}$ be a countable partition of $\Omega$ into events of positive probability, and let $\mathcal{G} = \sigma(\{B_n\})$ be the $\sigma$-algebra generated by the partition. Then
\[
E[X \mid \mathcal{G}](\omega) = \sum_n \frac{E[X \mathbf{1}_{B_n}]}{P(B_n)}\, \mathbf{1}_{B_n}(\omega) = \sum_n E[X \mid B_n]\, \mathbf{1}_{B_n}(\omega).
\]
In particular, taking $X = \mathbf{1}_A$ recovers the conditional probability: $E[\mathbf{1}_A \mid \mathcal{G}] = \sum_n P(A \mid B_n)\, \mathbf{1}_{B_n}$.
\end{example}

\begin{remark}[Measure-theoretic Bayes' rule]\label{rmk:bayes-measure}
Bayes' rule (\cref{thm:bayes}) extends to the measure-theoretic setting as follows. Let $P$ and $Q$ be two probability measures on $(\Omega, \cF)$ with $Q \ll P$, and let $L = dQ/dP$ be the Radon--Nikod\'ym derivative (the \emph{likelihood ratio}). Then for any sub-$\sigma$-algebra $\mathcal{G} \subseteq \cF$ and any bounded $\cF$-measurable $X$,
\begin{equation}\label{eq:abstract-bayes}
E_Q[X \mid \mathcal{G}] = \frac{E_P[XL \mid \mathcal{G}]}{E_P[L \mid \mathcal{G}]} \qquad Q\text{-a.s.}
\end{equation}
This is the abstract Bayes formula: it relates the conditional expectation under one measure to that under another, mediated by the likelihood ratio. The elementary Bayes' rule \eqref{eq:bayes} is the special case where $\mathcal{G} = \sigma(\{B_n\})$ and $Q(A) = P(A \mid B)$ for some fixed $B$.
\end{remark}

\begin{proposition}[Optimality of conditional expectation]\label{prop:cond-exp-optimal}
Let $Y \in L^2(\Omega, \cF, P)$ and $\mathcal{G} \subseteq \cF$ a sub-$\sigma$-algebra. Among all $\mathcal{G}$-measurable predictors $f \in L^2(\Omega, \mathcal{G}, P)$, the conditional expectation $E[Y \mid \mathcal{G}]$ uniquely minimizes the mean squared error:
\begin{equation}\label{eq:mse-decomp}
E[(Y - f)^2] = E[(Y - E[Y \mid \mathcal{G}])^2] + E[(E[Y \mid \mathcal{G}] - f)^2].
\end{equation}
\end{proposition}

\begin{proof}
Write $Y - f = (Y - E[Y \mid \mathcal{G}]) + (E[Y \mid \mathcal{G}] - f)$. The first term $\varepsilon = Y - E[Y \mid \mathcal{G}]$ is orthogonal to every element of $L^2(\Omega, \mathcal{G}, P)$ by the defining property of conditional expectation: for any $\mathcal{G}$-measurable $Z \in L^2$,
\[
E[\varepsilon \, Z] = E[(Y - E[Y \mid \mathcal{G}]) \, Z] = E[YZ] - E[E[Y \mid \mathcal{G}] \, Z] = 0.
\]
The second term $E[Y \mid \mathcal{G}] - f$ is $\mathcal{G}$-measurable, hence in $L^2(\Omega, \mathcal{G}, P)$. By Pythagoras in the Hilbert space $L^2(P)$:
\[
\|Y - f\|_2^2 = \|\varepsilon\|_2^2 + \|E[Y \mid \mathcal{G}] - f\|_2^2,
\]
which is \eqref{eq:mse-decomp}. The second term vanishes if and only if $f = E[Y \mid \mathcal{G}]$ a.s., giving uniqueness.
\end{proof}

\begin{corollary}[Monotonicity]\label{cor:monotonicity}
If $\mathcal{G}_1 \subseteq \mathcal{G}_2 \subseteq \cF$, then $E[(Y - E[Y \mid \mathcal{G}_1])^2] \geq E[(Y - E[Y \mid \mathcal{G}_2])^2]$. Finer $\sigma$-algebra $\to$ lower irreducible error. This is the mathematical content of the \emph{policy improvement theorem}: finer $\sigma$-algebra $\to$ lower projection residual.
\end{corollary}

\begin{proof}
Since $\mathcal{G}_1 \subseteq \mathcal{G}_2$, we have $L^2(\Omega, \mathcal{G}_1, P) \subseteq L^2(\Omega, \mathcal{G}_2, P)$. Projecting onto a larger closed subspace can only decrease the residual.
\end{proof}

\begin{definition}[The GPI category]\label{def:gpi-category}
Let $(\Omega, \cF, P)$ be a probability space and $Y \in L^2(\Omega, \cF, P)$.
The \emph{GPI category} $\Proj_Y$ is the thin category (poset category) whose
objects are sub-$\sigma$-algebras $\mathcal{G} \subseteq \cF$ and whose unique
morphism $\mathcal{G}_1 \to \mathcal{G}_2$ exists if and only if
$\mathcal{G}_1 \subseteq \mathcal{G}_2$.

The \emph{residual functor} is the contravariant functor
\[
R_Y : \Proj_Y^{\mathrm{op}} \to (\R_{\geq 0}, \geq), \qquad
R_Y(\mathcal{G}) = \|Y - E[Y \mid \mathcal{G}]\|_2^2.
\]
By \cref{cor:monotonicity}, $R_Y$ is a functor: every morphism
$\mathcal{G}_1 \hookrightarrow \mathcal{G}_2$ in $\Proj_Y$ maps to the
inequality $R_Y(\mathcal{G}_1) \geq R_Y(\mathcal{G}_2)$ in
$(\R_{\geq 0}, \geq)$.

The \emph{evaluation functor} is
\[
\mathcal{E}_Y : \Proj_Y \to L^2(\Omega, \cF, P), \qquad
\mathcal{E}_Y(\mathcal{G}) = E[Y \mid \mathcal{G}].
\]
A \emph{filtration} $\mathbb{F} = (\mathcal{G}_t)_{t \geq 0}$ with
$\mathcal{G}_0 \subseteq \mathcal{G}_1 \subseteq \cdots$ is a functor
$\mathbb{F} : (\mathbb{N}, \leq) \to \Proj_Y$ from the poset category of
natural numbers.
\end{definition}

\begin{example}[Generalized policy iteration and the $\sigma$-algebra lattice]\label{ex:ml-sigma}
We show that the central constructions of reinforcement learning and statistical learning are operations in the GPI category $\Proj_Y$ (\cref{def:gpi-category}). The core objects are the evaluation functor $\mathcal{E}_Y$ (\cref{def:gpi-category}) and an improvement operator $\mathcal{I}$, whose alternation is generalized policy iteration (GPI), a residual-monotone walk on $\Proj_Y$ governed by \cref{prop:cond-exp-optimal}.

\begin{enumerate}[label=(\alph*)]

\item \textbf{The two operators.} Let $Y \in L^2(\Omega, \cF, P)$ be a target random variable (a reward, a label, a future state). Let $\Lambda \subseteq \{\mathcal{G} : \mathcal{G} \subseteq \cF\}$ be the family of sub-$\sigma$-algebras reachable by the agent or model.

\emph{Evaluation operator} $\mathcal{E}$: given a fixed $\mathcal{G} \in \Lambda$, compute the orthogonal projection
\[
\mathcal{E}_{\mathcal{G}}[Y] = E[Y \mid \mathcal{G}].
\]
By \cref{prop:cond-exp-optimal}, this is the unique minimizer of $E[(Y - f)^2]$ over $f \in L^2(\Omega, \mathcal{G}, P)$. The $\sigma$-algebra is fixed; the function moves.

\emph{Improvement operator} $\mathcal{I}$: given the current projection residual, select a new $\sigma$-algebra $\mathcal{G}' \in \Lambda$ with smaller residual:
\[
\mathcal{I}(\mathcal{G}) = \mathcal{G}' \in \Lambda \quad \text{such that} \quad E[(Y - E[Y \mid \mathcal{G}'])^2] \leq E[(Y - E[Y \mid \mathcal{G}])^2].
\]
The target is fixed; the $\sigma$-algebra moves.

Note that $\mathcal{I}$ selects a $\sigma$-algebra with smaller \emph{residual}; this does not require $\mathcal{G} \subseteq \mathcal{G}'$ (the new $\sigma$-algebra need not be finer). When the walk is additionally a filtration ($\mathcal{G}_n \subseteq \mathcal{G}_{n+1}$ for all $n$), the stronger convergence of \cref{thm:in-context-gpi} applies---this is the case for in-context learning, where each new token extends the context.

\item \textbf{Generalized policy iteration.} GPI alternates the two operators:
\[
\mathcal{G}_0 \xrightarrow{\;\mathcal{E}\;} V_0 \xrightarrow{\;\mathcal{I}\;} \mathcal{G}_1 \xrightarrow{\;\mathcal{E}\;} V_1 \xrightarrow{\;\mathcal{I}\;} \mathcal{G}_2 \xrightarrow{\;\mathcal{E}\;} \cdots
\]
where $V_n = E[Y \mid \mathcal{G}_n]$. By \cref{cor:monotonicity}, each improvement step yields
\[
\|Y - V_{n+1}\|_2^2 \leq \|Y - V_n\|_2^2.
\]
The sequence of projection residuals is non-increasing and bounded below by zero, hence converges. In the finite-state setting, this convergence is achieved in finitely many steps and yields the optimal value function $V^* = E[Y \mid \mathcal{G}^*]$. This is the measure-theoretic skeleton of the \emph{policy improvement theorem}: the monotonicity of orthogonal projections on the lattice is the mechanism underlying the classical Bellman convergence argument.

\item \textbf{Reinforcement learning: the Bellman equation as tower property.} In a Markov decision process $(\mathcal{S}, \mathcal{A}, P, R, \gamma)$, the agent selects actions $A_t$ in states $S_t$, receiving rewards $R_{t+1}$. A policy $\pi$ determines the agent's behavior; the Markov property ensures
that the value function depends only on the current state, so the relevant
$\sigma$-algebra is $\sigma(S_t)$, a coarsening of the full trajectory
$\sigma$-algebra. The \emph{value function} is
\[
V^\pi(s) = E^\pi\!\left[\sum_{k=0}^\infty \gamma^k R_{t+k+1} \;\middle|\; S_t = s\right] = E[\mathcal{R}_t \mid \sigma(S_t)],
\]
where $\mathcal{R}_t = \sum_{k=0}^\infty \gamma^k R_{t+k+1}$ is the discounted return and $\gamma \in [0,1)$. This is the evaluation operator $\mathcal{E}$ applied to the return.

The \emph{Bellman equation} is the tower property in disguise:
\begin{equation}\label{eq:bellman}
V^\pi(s) = E^\pi[R_{t+1} + \gamma\, V^\pi(S_{t+1}) \mid S_t = s].
\end{equation}
Writing $\mathcal{G}_t = \sigma(S_t)$, this reads $E[\mathcal{R}_t \mid \mathcal{G}_t] = E\bigl[R_{t+1} + \gamma\, E[\mathcal{R}_{t+1} \mid \mathcal{G}_{t+1}] \mid \mathcal{G}_t\bigr]$---the tower property applied to the return, with the Markov property ensuring $\mathcal{G}_t$ captures the relevant information.

Policy improvement selects $\pi'$ greedy with respect to $V^\pi$: choose $A_t$ to maximize $E[R_{t+1} + \gamma V^\pi(S_{t+1}) \mid S_t, A_t]$. This generates a $\sigma$-algebra $\mathcal{G}_{\pi'}$ that resolves the action-value distinctions $\pi$ left unresolved, and by the same monotonicity principle, $V^{\pi'} \geq V^\pi$
(the classical policy improvement theorem; the pointwise inequality
follows from applying the Bellman operator's monotonicity state by state).

\item \textbf{Supervised learning as single-step evaluation.} The static supervised setting is the degenerate case of GPI: no dynamics, no sequential decisions. The joint distribution $(X, Y) \sim P$ is fixed. A model with parameters $\theta$ computes a representation $h_\theta : \mathcal{X} \to \R^d$ generating
\[
\mathcal{G}_\theta = \sigma(h_\theta(X)) = \{h_\theta^{-1}(B) : B \in \mathcal{B}(\R^d)\}.
\]
Two inputs $x, x'$ with $h_\theta(x) = h_\theta(x')$ are indistinguishable to the model---they belong to the same $\mathcal{G}_\theta$-atom, exactly as the Louvre and Prado Giocondas belong to the same $\mathcal{H}$-atom at 960px (\cref{ex:gioconda}). The evaluation operator gives the optimal readout $E[Y \mid \mathcal{G}_\theta]$; the improvement operator is stochastic gradient descent on $\theta$:
\begin{equation}\label{eq:ml-fundamental}
\min_\theta\, E[(Y - E[Y \mid \mathcal{G}_\theta])^2].
\end{equation}
There is no alternation---$\mathcal{E}$ and $\mathcal{I}$ are fused into a single optimization. GPI degenerates to gradient descent on the projection residual.

\item \textbf{Architecture $=$ reachable $\sigma$-algebras.} Different architectures constrain which $\sigma$-algebras are reachable:
\begin{itemize}
\item A \emph{linear model}: $\mathcal{G}_\theta = \sigma(w^\top X)$---a hyperplane partition. Coarse.
\item A \emph{CNN}: $\mathcal{G}_\theta$ respects translation equivariance---only translation-invariant distinctions are visible.
\item A \emph{Transformer} with context length $T$: $\mathcal{G}_\theta = \sigma(h_\theta(x_1, \ldots, x_T))$---the $\sigma$-algebra generated by attending to all positions in the context window.
\end{itemize}
Architecture design is the art of choosing which region of the lattice $\Lambda$ to search over.

\item \textbf{Overfitting, underfitting, and regularization.}
\begin{itemize}
\item \emph{Underfitting}: $\mathcal{G}_\theta$ is too coarse---important distinctions are invisible. The projection $E[Y \mid \mathcal{G}_\theta]$ has large residual.
\item \emph{Overfitting}: $\mathcal{G}_\theta$ is too fine on training data---it separates noise from noise. On test data from the same $P$, these spurious distinctions do not hold.
\item \emph{Regularization} (dropout, weight decay, early stopping) $=$ constraining $\Lambda$, preventing $\mathcal{G}_\theta$ from becoming too fine relative to the sample size.
\end{itemize}

\item \textbf{Next-token prediction (GPT).} The sequence $(X_1, X_2, \ldots)$ is drawn from $P$. The autoregressive model computes
\[
f_\theta(x_1, \ldots, x_t) \approx E[X_{t+1} \mid \sigma(X_1, \ldots, X_t)]
\]
for each $t$. This is GPI with a rolling evaluation: at each position, $\mathcal{E}$ projects the next token onto the $\sigma$-algebra generated by the context, and $\mathcal{I}$ updates $\theta$ to reduce the cross-entropy (the KL divergence from the model's conditional to the true conditional). The transformer's attention mechanism parameterizes the map from context to the $\sigma$-algebra $\mathcal{G}_\theta^{(t)}$ that determines what the model ``remembers.''

\item \textbf{World models (JEPA).} A world model predicts future states given current observations. In the framework:
\[
\text{JEPA predictor} = E[h_\theta(X_{t+1}) \mid \mathcal{G}_\theta^{(t)}].
\]
The key design choice: predict in \emph{representation space} $h_\theta(X_{t+1})$, not in pixel space $X_{t+1}$. This means the target itself lives in a learned $\sigma$-algebra. Both $\mathcal{E}$ and its target co-evolve---the model is simultaneously choosing what to predict and how to predict it. This is joint optimization over two $\sigma$-algebras: GPI where the improvement operator $\mathcal{I}$ moves both the predictor and the target.

\item \textbf{Diffusion models.} The forward process $X_0 \to X_1 \to \cdots \to X_T$ is a Markov chain that progressively destroys information: the mutual information $I(X_0; X_t)$ decreases with $t$, so the conditional expectation $E[f(X_0) \mid X_t]$ becomes coarser as $t$ grows. At $t = T$, the conditional distribution is nearly independent of $X_0$ (pure noise). The reverse (generative) process recovers information step by step---this is GPI run backwards on the lattice. The score function $\nabla_{x_t} \log p(x_t)$ is the gradient of the log-density, and score matching learns to invert each information-destroying step. Generation $=$ climbing from uninformative to informative, one evaluation--improvement pair at a time.

\end{enumerate}
\end{example}

\begin{remark}[The GPI--measure-theory dictionary]\label{rmk:ml-dictionary}
The following table summarizes the translation developed in \cref{ex:ml-sigma}.

\medskip
\begin{center}
\renewcommand{\arraystretch}{1.25}
\small
\begin{tabular}{@{}p{0.38\textwidth}p{0.57\textwidth}@{}}
\hline
\textbf{AI concept} & \textbf{Measure-theoretic object} \\
\hline
Target (reward, label, future state) $Y$ & Random variable in $L^2(\Omega, \cF, P)$ \\
Evaluation operator $\mathcal{E}$ & Orthogonal projection $E[Y \mid \mathcal{G}]$ \\
Improvement operator $\mathcal{I}$ & Lattice navigation $\mathcal{G} \mapsto \mathcal{G}'$ \\
GPI convergence & Monotone walk on $\Lambda$, bounded $\Rightarrow$ converges \\
Bellman equation & Tower property of conditional expectation \\
Policy improvement theorem & \cref{cor:monotonicity}: finer $\mathcal{G} \to$ smaller residual \\
Value function $V^\pi$ & $E[\mathcal{R}_t \mid \sigma(S_t)]$ (projection of return) \\
\hline
Learned representation $h_\theta(X)$ & Generator of $\mathcal{G}_\theta = \sigma(h_\theta(X))$ \\
Architecture constraints & Reachable sub-lattice $\Lambda \subseteq \{\mathcal{G} \subseteq \cF\}$ \\
Supervised training loss & $E[(Y - E[Y \mid \mathcal{G}_\theta])^2]$ (projection residual) \\
Training over $\theta$ & Search over $\sigma$-algebras in $\Lambda$ \\
Overfitting & $\mathcal{G}_\theta$ too fine (separates noise) \\
Regularization & Constraining $\Lambda$ \\
\hline
Next-token prediction & $E[X_{t+1} \mid \sigma(X_1, \ldots, X_t)]$ \\
World model (JEPA) & $E[h_\theta(X_{t+1}) \mid \mathcal{G}_\theta^{(t)}]$ (joint $\sigma$-algebra opt.) \\
Diffusion forward & Information destruction: $I(X_0; X_T) \ll I(X_0; X_0)$ \\
Diffusion reverse & Information recovery: climbing toward finer conditional structure \\
Generalization & $\mathcal{G}_\theta$ captures structure of $P$, not artifacts \\
\hline
GPI category $\Proj_Y$ & Thin category $(\Lambda, \subseteq)$ (\cref{def:gpi-category}) \\
In-context learning & Filtration $\mathbb{F} : (\mathbb{N}, \leq) \to \Proj_Y$ \\
In-context convergence & Colimit $\mathcal{G}_\infty$ in $\Proj_Y$ (\cref{thm:in-context-gpi}) \\
\hline
\end{tabular}
\normalsize
\end{center}
\medskip

\noindent Every model, every architecture, every training procedure in modern AI is a particular strategy for navigating the lattice of sub-$\sigma$-algebras of $\cF$. Reinforcement learning makes the navigation explicit through the alternation $\mathcal{E} \circ \mathcal{I}$; supervised learning fuses the two into a single gradient step. The lattice itself---and the projection theorem that governs it---is Kolmogorov's, 1933.

\smallskip
\noindent\textit{Author's note.} Machine learning, as practiced, does not resemble the popular image of ``artificial intelligence''---it resembles $\sigma$-algebra selection. The objects manipulated by gradient descent are not thoughts, beliefs, or goals; they are partitions of a sample space. The questions answered by a trained model are not ``what should I do?''\ but ``which distinctions does the data support?''\ The $\sigma$-algebra lattice is a more faithful portrait of what these systems are and do than any anthropomorphic metaphor.
\end{remark}

\begin{theorem}[In-context GPI]\label{thm:in-context-gpi}
Let $Y \in L^2(\Omega, \cF, P)$ and let
$\mathbb{F} : (\mathbb{N}, \leq) \to \Proj_Y$ be a filtration, with
$\mathbb{F}(t) = \mathcal{G}_t$ and
$\mathcal{G}_0 \subseteq \mathcal{G}_1 \subseteq \cdots \subseteq \cF$.
\begin{enumerate}[label=(\roman*)]
\item \emph{(Monotone residual.)} The composite
$R_Y \circ \mathbb{F}^{\mathrm{op}}$ is order-preserving:
$R_Y(\mathcal{G}_t)$ is non-increasing in $t$.

\item \emph{(Colimit convergence.)} Let
$\mathcal{G}_\infty = \sigma\!\bigl(\bigcup_{t \geq 0} \mathcal{G}_t\bigr)$
be the colimit of $\mathbb{F}$ in $\Proj_Y$. Then
\[
\mathcal{E}_Y(\mathcal{G}_t) = E[Y \mid \mathcal{G}_t]
\xrightarrow{\;L^2\;} E[Y \mid \mathcal{G}_\infty]
= \mathcal{E}_Y(\mathcal{G}_\infty)
\quad \text{as } t \to \infty.
\]
The evaluation functor sends the colimit of the filtration to the
$L^2$-limit of the projections.

\item \emph{(Terminal projection.)}
$R_Y(\mathcal{G}_\infty) = \inf_{t \geq 0} R_Y(\mathcal{G}_t)$.
The colimit achieves the infimum of the residual functor over the image
of $\mathbb{F}$.
\end{enumerate}
\end{theorem}

\begin{proof}
(i) is \cref{cor:monotonicity} applied to each inclusion
$\mathcal{G}_t \subseteq \mathcal{G}_{t+1}$.

(ii) The sequence $V_t = E[Y \mid \mathcal{G}_t]$ is a martingale
adapted to $(\mathcal{G}_t)$: for $s \leq t$,
$E[V_t \mid \mathcal{G}_s] = E[Y \mid \mathcal{G}_s] = V_s$ by the tower
property. Since $\sup_t E[V_t^2] \leq E[Y^2] < \infty$, L\'{e}vy's upward
theorem \cite[Theorem~4.4.5]{Durrett} gives
$V_t \to E[Y \mid \mathcal{G}_\infty]$ both a.s.\ and in $L^2$.

(iii) By (ii) and continuity of the $L^2$-norm:
$R_Y(\mathcal{G}_t) = \|Y - V_t\|_2^2 \to
\|Y - E[Y \mid \mathcal{G}_\infty]\|_2^2 = R_Y(\mathcal{G}_\infty)$.
Combined with (i), $R_Y(\mathcal{G}_\infty)$ is the infimum.
\end{proof}

\begin{remark}[In-context reinforcement learning]\label{rmk:in-context-rl}
\cref{thm:in-context-gpi} gives the measure-theoretic content of
\emph{in-context reinforcement learning}: a transformer processing a
sequence of $(s_k, a_k, r_{k+1})$ tuples builds a filtration
$\mathcal{G}_t = \sigma(s_1, a_1, r_2, \ldots, s_t)$ within its context
window. Each new token is a morphism
$\mathcal{G}_t \hookrightarrow \mathcal{G}_{t+1}$ in $\Proj_Y$. The three
modes of GPI are distinguished by the role of the improvement operator:

\medskip
\begin{center}
\renewcommand{\arraystretch}{1.25}
\small
\begin{tabular}{@{}p{0.22\textwidth}p{0.34\textwidth}p{0.36\textwidth}@{}}
\hline
\textbf{Mode} & \textbf{Improvement $\mathcal{I}$} & \textbf{Categorical structure} \\
\hline
Classical RL & Explicit: $\pi \to \pi'$ & Endofunctor on $\Proj_Y$ \\
Supervised learning & Fused: $\nabla_\theta$ on $R_Y(\mathcal{G}_\theta)$ & Search over $\Lambda \subset \mathrm{Ob}(\Proj_Y)$ \\
In-context RL & Implicit: data extends context & Filtration $\mathbb{F} : (\mathbb{N}, \leq) \to \Proj_Y$ \\
\hline
\end{tabular}
\normalsize
\end{center}
\medskip

\noindent In the in-context case, the improvement operator is the identity
on the filtration---improvement is free, supplied by the data. The
transformer's forward pass at position $t$ approximates the evaluation
functor $\mathcal{E}_Y(\mathcal{G}_t)$, and convergence is guaranteed by
\cref{thm:in-context-gpi}(ii).
\end{remark}

\begin{remark}[Connection to the CLT framework]
The foundations reviewed here underlie every result in this paper. The CLT (\cref{thm:add-clt}) concerns the law $\mu_{S_n/\sqrt{n}} = P \circ (S_n/\sqrt{n})^{-1}$ of a normalized sum, whose characteristic function is an expectation: $\varphi_{S_n/\sqrt{n}}(\xi) = E[e^{i\xi S_n/\sqrt{n}}]$. The Fourier transform of a probability measure is nothing but the expectation of a character. The entire duality between the additive and multiplicative CLTs (\cref{thm:mellin-fourier}) rests on the fact that expectation is a linear functional that commutes with the group isomorphism $\log : (\Rpos, \times) \to (\R, +)$.
\end{remark}

%% ============================================================
\section{Preliminaries on locally compact abelian groups}

We recall the basic setup of harmonic analysis on locally compact abelian groups. Standard references include Rudin \cite{Rudin}, Folland \cite{Folland}, and Hewitt--Ross \cite{HewittRoss}.

\begin{definition}\label{def:LCA}
A \emph{locally compact abelian (LCA) group} is a topological group $G$ that is abelian, Hausdorff, and locally compact. By the Haar theorem \cite[Theorem~2.10]{Folland}, there exists a Radon measure $\mu_G$ on $G$, unique up to positive scalar multiple, that is invariant under translation. This measure is called the \emph{Haar measure} on $G$.
\end{definition}

\begin{definition}\label{def:dual}
A \emph{character} of an LCA group $G$ is a continuous group homomorphism $\chi : G \to \bT$, where $\bT = \{z \in \mathbb{C} : |z| = 1\}$ is the circle group under multiplication. The set of all characters forms an LCA group $\widehat{G}$ under pointwise multiplication, called the \emph{Pontryagin dual} of $G$.
\end{definition}

\begin{example}\label{ex:groups}
The following table summarizes our three running examples.
\begin{center}
\renewcommand{\arraystretch}{1.4}
\begin{tabular}{lccc}
\textbf{Group $G$} & \textbf{Haar measure} & \textbf{Dual $\widehat{G}$} & \textbf{Characters} \\
\hline
$(\R, +)$ & $dx$ & $\R$ & $\chi_\xi(x) = e^{i\xi x}$ \\
$(\Rpos, \times)$ & $dx/x$ & $\R$ & $\chi_s(x) = x^{is}$ \\
$(\bT, +)$ & $d\theta$ & $\Z$ & $\chi_n(\theta) = e^{2\pi i n\theta}$
\end{tabular}
\end{center}
Note that $(\R, +)$ and $(\Rpos, \times)$ have isomorphic duals but are distinct as groups.
\end{example}

\begin{definition}\label{def:FT}
Let $\mu$ be a probability measure on an LCA group $G$. The \emph{Fourier transform} (or \emph{characteristic function}) of $\mu$ is the function $\widehat{\mu} : \widehat{G} \to \mathbb{C}$ defined by
\[
\widehat{\mu}(\chi) = \int_G \chi(x) \, d\mu(x).
\]
\end{definition}

When $G = (\R, +)$, this reduces to the usual characteristic function $\widehat{\mu}(\xi) = \int_\R e^{i\xi x}\, d\mu(x)$. When $G = (\Rpos, \times)$, this gives $\widehat{\mu}(s) = \int_0^\infty x^{is}\, d\mu(x)$, which we call the Mellin transform.

\begin{theorem}[Pontryagin Duality {\cite[Theorem~1.7.2]{Rudin}}]\label{thm:pontryagin}
For any LCA group $G$, the canonical map $G \to \widehat{\widehat{G}}$ defined by $x \mapsto (\chi \mapsto \chi(x))$ is an isomorphism of topological groups.
\end{theorem}

\begin{theorem}[L\'evy Continuity Theorem on LCA Groups {\cite[Theorem~VI.4.2]{Parthasarathy}}]\label{thm:levy}
Let $\{\mu_n\}$ be a sequence of probability measures on an LCA group $G$, and let $\mu$ be a probability measure on $G$. Then $\mu_n \xrightarrow{w} \mu$ if and only if $\widehat{\mu_n}(\chi) \to \widehat{\mu}(\chi)$ for every $\chi \in \widehat{G}$.
\end{theorem}

%% ============================================================
\section{The additive central limit theorem}

We work on the group $G = (\R, +)$ with Haar measure $dx$, dual $\widehat{G} = \R$, and characters $\chi_\xi(x) = e^{i\xi x}$.

\begin{theorem}[Classical CLT]\label{thm:add-clt}
Let $X_1, X_2, \ldots$ be i.i.d.\ real-valued random variables with $E[X_1] = 0$ and $\Var(X_1) = \sigma^2 \in (0, \infty)$. Define $S_n = \sum_{k=1}^n X_k$. Then
\[
\frac{S_n}{\sqrt{n}} \xrightarrow{d} \cN(0, \sigma^2).
\]
\end{theorem}

\begin{proof}
Let $\varphi(\xi) = E[e^{i\xi X_1}]$ denote the characteristic function of $X_1$. Since $E[X_1] = 0$ and $E[X_1^2] = \sigma^2 < \infty$, we have the expansion
\[
\varphi(\xi) = 1 - \frac{\sigma^2 \xi^2}{2} + o(\xi^2) \quad \text{as } \xi \to 0.
\]
By independence, the characteristic function of $S_n / \sqrt{n}$ is
\[
\varphi_{S_n/\sqrt{n}}(\xi)
= \left[\varphi\!\left(\frac{\xi}{\sqrt{n}}\right)\right]^n
= \left[1 - \frac{\sigma^2 \xi^2}{2n} + o\!\left(\frac{1}{n}\right)\right]^n.
\]
Taking $n \to \infty$:
\[
\left[1 - \frac{\sigma^2 \xi^2}{2n} + o\!\left(\frac{1}{n}\right)\right]^n \longrightarrow e^{-\sigma^2 \xi^2 / 2},
\]
which is the characteristic function of $\cN(0, \sigma^2)$. The conclusion follows from the L\'evy continuity theorem (\cref{thm:levy}).
\end{proof}

\begin{remark}\label{rmk:gaussian-dual}
The characteristic function $\xi \mapsto e^{-\sigma^2\xi^2/2}$ is itself a Gaussian function on the dual group $\widehat{\R} \cong \R$. This self-similarity under the Fourier transform is not accidental; it reflects the fact that the Gaussian is an eigenfunction of $\cF$ (see \cref{thm:gaussian-fixed}).
\end{remark}

\begin{proposition}[Heat equation]\label{prop:heat-add}
The Gaussian density $p_t(x) = (2\pi t)^{-1/2}\, e^{-x^2/(2t)}$ is the fundamental solution of the heat equation
\begin{equation}\label{eq:heat-add}
\partial_t u = \tfrac{1}{2}\, \partial_{xx} u
\end{equation}
on $(\R, +)$. In Fourier space, \eqref{eq:heat-add} becomes the ODE
\begin{equation}\label{eq:heat-add-fourier}
\partial_t \widehat{u}(\xi) = -\tfrac{\xi^2}{2}\, \widehat{u}(\xi),
\end{equation}
with solution $\widehat{u}(t, \xi) = e^{-t\xi^2/2}$. The family $\{p_t\}_{t > 0}$ forms a convolution semigroup: $p_t * p_s = p_{t+s}$.
\end{proposition}

\begin{proof}
Apply the Fourier transform to \eqref{eq:heat-add}. Since $\widehat{\partial_{xx} u}(\xi) = (i\xi)^2 \widehat{u}(\xi) = -\xi^2 \widehat{u}(\xi)$, the PDE reduces to the ODE \eqref{eq:heat-add-fourier}, whose unique solution with initial condition $\widehat{u}(0, \xi) = 1$ is $\widehat{u}(t, \xi) = e^{-t\xi^2/2}$. Inverting the Fourier transform yields $p_t$. The semigroup property follows from $e^{-t\xi^2/2} \cdot e^{-s\xi^2/2} = e^{-(t+s)\xi^2/2}$.
\end{proof}

%% ============================================================
\section{The multiplicative central limit theorem}

We now work on the group $G = (\Rpos, \times)$ with Haar measure $dx/x$, dual $\widehat{G} = \R$, and characters $\chi_s(x) = x^{is}$.

\begin{definition}\label{def:mellin}
The \emph{Mellin transform} of a probability measure $\mu$ on $\Rpos$ is
\[
\cM[\mu](s) = \int_0^\infty x^{is}\, d\mu(x), \qquad s \in \R.
\]
This is the Fourier transform of $\mu$ viewed as a measure on the LCA group $(\Rpos, \times)$.
\end{definition}

\begin{theorem}[Multiplicative CLT]\label{thm:mult-clt}
Let $Y_1, Y_2, \ldots$ be i.i.d.\ positive random variables with $E[\log Y_1] = 0$ and $\Var(\log Y_1) = \sigma^2 \in (0, \infty)$. Define $P_n = \prod_{k=1}^n Y_k$. Then
\[
P_n^{1/\sqrt{n}} \xrightarrow{d} \operatorname{LogNormal}(0, \sigma^2).
\]
\end{theorem}

\begin{proof}
Set $X_k = \log Y_k$. Then $X_1, X_2, \ldots$ are i.i.d.\ with $E[X_k] = 0$ and $\Var(X_k) = \sigma^2$. Since
\[
\log\!\left(P_n^{1/\sqrt{n}}\right) = \frac{1}{\sqrt{n}} \sum_{k=1}^n \log Y_k = \frac{S_n}{\sqrt{n}},
\]
\cref{thm:add-clt} gives $S_n/\sqrt{n} \xrightarrow{d} \cN(0, \sigma^2)$. As $\exp : \R \to \Rpos$ is continuous, the continuous mapping theorem yields
\[
P_n^{1/\sqrt{n}} = \exp\!\left(\frac{S_n}{\sqrt{n}}\right) \xrightarrow{d} \exp\!\left(\cN(0, \sigma^2)\right) = \operatorname{LogNormal}(0, \sigma^2). \qedhere
\]
\end{proof}

\begin{remark}[Normalization]\label{rmk:normalization}
On $(\R, +)$, the natural scaling is $x \mapsto x/\sqrt{n}$ (division by $\sqrt{n}$). On $(\Rpos, \times)$, the corresponding operation is $y \mapsto y^{1/\sqrt{n}}$ (extraction of the $\sqrt{n}$-th root), since $\log(y^{1/\sqrt{n}}) = (\log y)/\sqrt{n}$. Division on the additive group corresponds to root extraction on the multiplicative group.
\end{remark}

\begin{proposition}[Multiplicative heat equation]\label{prop:heat-mult}
The log-normal density (with respect to Lebesgue measure on $\Rpos$)
\[
q_t(x) = \frac{1}{x\sqrt{2\pi t}}\, \exp\!\left(-\frac{(\log x)^2}{2t}\right), \qquad x > 0,
\]
is the fundamental solution of the heat equation on $(\Rpos, \times)$:
\begin{equation}\label{eq:heat-mult}
\partial_t u = \tfrac{1}{2}\, D^2 u, \qquad D = x\partial_x.
\end{equation}
Here $D = x\partial_x$ is the Euler operator, which is the left-invariant vector field on $(\Rpos, \times)$.
\end{proposition}

\begin{proof}
Set $v(t, y) = u(t, e^y)$. Then $Du(t,x)\big|_{x=e^y} = \partial_y v(t,y)$, and consequently $D^2 u\big|_{x=e^y} = \partial_{yy}v$. Thus \eqref{eq:heat-mult} transforms into the standard heat equation $\partial_t v = \frac{1}{2}\partial_{yy}v$ on $\R$, whose fundamental solution is $p_t(y)$. Substituting $y = \log x$ yields $u(t,x) = q_t(x)$.
\end{proof}

\begin{remark}\label{rmk:euler-invariant}
The Euler operator $D = x\partial_x$ is left-invariant on $(\Rpos, \times)$: for any $a > 0$ and smooth $f$,
\[
D\big[f(a\,\cdot\,)\big](x) = x\,\frac{d}{dx}f(ax) = ax\,f'(ax) = (Df)(ax),
\]
so $D \circ L_a^* = L_a^* \circ D$, where $L_a$ denotes left multiplication by $a$. The operator $D$ generates the one-parameter group of ``dilations'' $x \mapsto x^{e^t} = e^{t\log x}$ and plays the same role on $(\Rpos, \times)$ that $\partial_x$ plays on $(\R, +)$.
\end{remark}

%% ============================================================
\section{The Fourier--Mellin duality}

We now establish the central result: the Fourier and Mellin transforms are related by the group isomorphism $\log : (\Rpos, \times) \to (\R, +)$.

\begin{proposition}\label{prop:log-iso}
The map $\log : (\Rpos, \times) \to (\R, +)$ is an isomorphism of LCA groups. It carries the Haar measure $dx/x$ on $\Rpos$ to the Lebesgue measure $dy$ on $\R$.
\end{proposition}

\begin{proof}
\emph{Group homomorphism.} For all $a, b > 0$: $\log(ab) = \log a + \log b$.

\emph{Topological isomorphism.} The map $\log : (0, \infty) \to \R$ is a homeomorphism with continuous inverse $\exp$.

\emph{Haar measure.} Under the substitution $y = \log x$ (i.e., $x = e^y$, $dx = e^y\, dy$):
\[
\frac{dx}{x} = \frac{e^y\, dy}{e^y} = dy.
\]
Thus $\log_*(dx/x) = dy$, confirming that $\log$ intertwines the respective Haar measures.
\end{proof}

\begin{theorem}[Mellin--Fourier correspondence]\label{thm:mellin-fourier}
Let $\mu$ be a probability measure on $(\Rpos, \times)$, and let $\log_*\mu$ denote its pushforward to $(\R, +)$ under $\log$. Then
\begin{equation}\label{eq:mellin-fourier}
\cM[\mu](s) = \cF[\log_*\mu](s) \qquad \text{for all } s \in \R.
\end{equation}
That is, the following diagram commutes:
\begin{equation}\label{eq:cd-measures}
\begin{tikzcd}
\cP(\Rpos, \times) \arrow[r, "\widehat{(\cdot)}_{\cM}"] \arrow[d, "\log_*"'] & C_b(\R) \arrow[d, "="] \\
\cP(\R, +) \arrow[r, "\widehat{(\cdot)}_{\cF}"'] & C_b(\R)
\end{tikzcd}
\end{equation}
where $\cP(\cdot)$ denotes the space of probability measures and $C_b(\R)$ the space of bounded continuous functions.
\end{theorem}

\begin{proof}
For any $s \in \R$,
\begin{align*}
\cM[\mu](s) &= \int_0^\infty x^{is}\, d\mu(x) \\
&= \int_{-\infty}^{\infty} e^{isy}\, d(\log_*\mu)(y) \qquad \text{(substituting } y = \log x,\; x^{is} = e^{is\log x}\text{)} \\
&= \cF[\log_*\mu](s). \qedhere
\end{align*}
\end{proof}

\begin{corollary}\label{cor:clt-image}
The multiplicative CLT (\cref{thm:mult-clt}) is the image of the additive CLT (\cref{thm:add-clt}) under the group isomorphism $\exp : (\R, +) \to (\Rpos, \times)$. Specifically, the log-normal distribution is the pushforward of the Gaussian under $\exp$:
\[
\operatorname{LogNormal}(0, \sigma^2) = \exp_*\, \cN(0, \sigma^2).
\]
\end{corollary}

\begin{remark}[Functoriality]\label{rmk:functoriality}
\cref{thm:mellin-fourier} is a special case of the functoriality of the Fourier transform on LCA groups. If $\phi : G_1 \to G_2$ is a continuous homomorphism of LCA groups and $\widehat{\phi} : \widehat{G_2} \to \widehat{G_1}$ is the induced dual map $\widehat{\phi}(\chi) = \chi \circ \phi$, then for any probability measure $\mu$ on $G_1$:
\[
\widehat{\phi_*\mu}(\chi) = \widehat{\mu}(\chi \circ \phi) = \widehat{\mu}\!\left(\widehat{\phi}(\chi)\right).
\]
In our setting, $\phi = \log$ and $\phi^{-1} = \exp$. The dual map $\widehat{\log} : \widehat{(\R,+)} \to \widehat{(\Rpos,\times)}$ sends the character $\chi_s$ of $(\R,+)$ (namely $y \mapsto e^{isy}$) to the character $\chi_s \circ \log$ of $(\Rpos,\times)$ (namely $x \mapsto e^{is\log x} = x^{is}$). Since both duals are identified with $\R$ via the parameter $s$, the dual map is the identity.
\end{remark}

\begin{proposition}[Duality of heat equations]\label{prop:heat-duality}
Under the isomorphism $\log : (\Rpos, \times) \to (\R, +)$, the multiplicative heat equation \eqref{eq:heat-mult} is conjugate to the additive heat equation \eqref{eq:heat-add}. The following diagram commutes:
\begin{equation}\label{eq:cd-heat}
\begin{tikzcd}[column sep=large, row sep=large]
\partial_t u = \tfrac{1}{2}D^2 u \;\text{on}\; \Rpos \arrow[r, leftrightarrow, "\cM"] \arrow[d, leftrightarrow, "u(t{,}x) {=} v(t{,}\log x)"'] & \partial_t \widehat{u} = -\tfrac{s^2}{2}\,\widehat{u} \arrow[d, "="] \\
\partial_t v = \tfrac{1}{2}\,\partial_{yy} v \;\text{on}\; \R \arrow[r, leftrightarrow, "\cF"'] & \partial_t \widehat{v} = -\tfrac{\xi^2}{2}\,\widehat{v}
\end{tikzcd}
\end{equation}
In particular, both heat equations reduce to the same ODE in the frequency domain.
\end{proposition}

\begin{proof}
The left vertical arrow is the content of \cref{prop:heat-mult}: the substitution $v(t,y) = u(t, e^y)$ conjugates $D^2 = (x\partial_x)^2$ to $\partial_{yy}$. The horizontal arrows are the Mellin and Fourier transforms, each reducing the respective PDE to the ODE $\partial_t \widehat{w} = -\frac{s^2}{2}\widehat{w}$. The right vertical arrow is the identity because, by \cref{thm:mellin-fourier}, the Mellin transform of $u$ in $x$ equals the Fourier transform of $v$ in $y$.
\end{proof}

%% ============================================================
\section{Gaussian measures on LCA groups}

We now place the Gaussian and log-normal distributions in the general framework of Gaussian measures on LCA groups, following Parthasarathy \cite{Parthasarathy} and Heyer \cite{Heyer}.

\begin{definition}[{cf.\ \cite[Definition~VI.3.1]{Parthasarathy}}]\label{def:quadratic}
A function $Q : \widehat{G} \to \R$ on the Pontryagin dual of an LCA group $G$ is a \emph{non-negative definite quadratic form} if:
\begin{enumerate}[label=(\roman*)]
\item $Q$ is continuous and $Q(\chi) \geq 0$ for all $\chi \in \widehat{G}$;
\item $Q$ satisfies the parallelogram law:
\[
Q(\chi_1 \chi_2) + Q(\chi_1 \chi_2^{-1}) = 2Q(\chi_1) + 2Q(\chi_2) \qquad \text{for all } \chi_1, \chi_2 \in \widehat{G};
\]
\item $Q(e) = 0$, where $e$ is the identity character.
\end{enumerate}
\end{definition}

\begin{definition}[{cf.\ \cite[Definition~VI.5.1]{Parthasarathy}}]\label{def:gaussian-lca}
A probability measure $\gamma$ on an LCA group $G$ is called \emph{Gaussian} if its Fourier transform has the form
\[
\widehat{\gamma}(\chi) = \chi(a) \cdot \exp\!\left(-Q(\chi)\right)
\]
for some $a \in G$ (the \emph{mean}) and some non-negative definite quadratic form $Q$ on $\widehat{G}$ (the \emph{covariance}).
\end{definition}

\begin{example}\label{ex:gaussian-instances}
We verify that \cref{def:gaussian-lca} recovers the classical examples.
\begin{enumerate}[label=(\alph*)]
\item \textbf{$G = (\R, +)$.} Here $\widehat{G} = \R$ with $\chi_\xi(x) = e^{i\xi x}$. Setting $a = \mu$ and $Q(\xi) = \sigma^2\xi^2/2$, we obtain $\widehat{\gamma}(\xi) = e^{i\mu\xi - \sigma^2\xi^2/2}$, the characteristic function of $\cN(\mu, \sigma^2)$. The parallelogram law is verified by direct computation:
\[
Q(\xi_1 + \xi_2) + Q(\xi_1 - \xi_2) = \tfrac{\sigma^2}{2}\bigl[(\xi_1+\xi_2)^2 + (\xi_1-\xi_2)^2\bigr] = \sigma^2(\xi_1^2 + \xi_2^2) = 2Q(\xi_1) + 2Q(\xi_2).
\]

\item \textbf{$G = (\Rpos, \times)$.} Here $\widehat{G} = \R$ with $\chi_s(x) = x^{is}$. Setting $a = e^\mu$ and $Q(s) = \sigma^2 s^2/2$ (the same quadratic form as in~(a)), we obtain $\widehat{\gamma}(s) = e^{i\mu s - \sigma^2 s^2/2}$, the Mellin transform of $\operatorname{LogNormal}(\mu, \sigma^2)$. By \cref{thm:mellin-fourier}, the quadratic forms on $\widehat{G} = \R$ are identical in both cases, reflecting the isomorphism of the dual groups.

\item \textbf{$G = (\bT, +)$.} Here $\widehat{G} = \Z$ with $\chi_n(\theta) = e^{2\pi i n\theta}$. Setting $Q(n) = \sigma^2 n^2/2$, we obtain $\widehat{\gamma}(n) = e^{-\sigma^2 n^2/2}$, the Fourier coefficients of the \emph{wrapped normal distribution}. Its density on $[0, 1)$ is
\[
f(\theta) = \sum_{k \in \Z} \frac{1}{\sqrt{2\pi}\,\sigma}\, \exp\!\left(-\frac{(\theta - k)^2}{2\sigma^2}\right) = \frac{1}{\sqrt{2\pi}\,\sigma}\,\vartheta_3\!\left(\pi\theta,\, e^{-\sigma^2/2}\right),
\]
where $\vartheta_3$ is the Jacobi theta function.
\end{enumerate}
\end{example}

\begin{theorem}[Gaussian as Fourier eigenfunction]\label{thm:gaussian-fixed}
Consider the Plancherel-normalized Fourier transform on $L^2(\R)$:
\[
\hat{f}(\xi) = \int_{-\infty}^{\infty} f(x)\, e^{-2\pi i x\xi}\, dx.
\]
The function $g(x) = e^{-\pi x^2}$ satisfies $\hat{g} = g$. That is, $g$ is an eigenfunction of $\cF$ with eigenvalue $1$.
\end{theorem}

\begin{proof}
Define $h(\xi) = \hat{g}(\xi) = \int_{-\infty}^{\infty} e^{-\pi x^2}\, e^{-2\pi i x\xi}\, dx$. Differentiating under the integral:
\[
h'(\xi) = -2\pi i \int_{-\infty}^{\infty} x\, e^{-\pi x^2}\, e^{-2\pi i x\xi}\, dx.
\]
Since $-2\pi x\, e^{-\pi x^2} = \frac{d}{dx}\!\left(e^{-\pi x^2}\right)$, we integrate by parts:
\begin{align*}
h'(\xi) &= i \int_{-\infty}^{\infty} \frac{d}{dx}\!\left(e^{-\pi x^2}\right) e^{-2\pi i x\xi}\, dx \\
&= -i \int_{-\infty}^{\infty} e^{-\pi x^2} \cdot (-2\pi i\xi)\, e^{-2\pi i x\xi}\, dx \\
&= -2\pi\xi \int_{-\infty}^{\infty} e^{-\pi x^2}\, e^{-2\pi i x\xi}\, dx \;=\; -2\pi\xi\, h(\xi).
\end{align*}
Thus $h$ satisfies the ODE $h' = -2\pi\xi\, h$ with initial condition $h(0) = \int_{-\infty}^{\infty} e^{-\pi x^2}\, dx = 1$. The unique solution is $h(\xi) = e^{-\pi\xi^2} = g(\xi)$.
\end{proof}

\begin{remark}\label{rmk:eigenvalues}
The eigenvalues of the Fourier transform on $L^2(\R)$ are $\{1, -1, i, -i\}$ (since $\cF^4 = \mathrm{id}$), with eigenspaces spanned by the Hermite functions $\psi_n(x) = H_n(\sqrt{2\pi}\, x)\, e^{-\pi x^2}$, where $H_n$ is the $n$-th (probabilist's) Hermite polynomial. Specifically, $\cF[\psi_n] = (-i)^n \psi_n$. The Gaussian $g = \psi_0$ lies in the eigenspace for eigenvalue $1$. This spectral decomposition is intimately connected to the Stone--von Neumann theorem and the metaplectic representation of $\mathrm{SL}(2, \R)$.
\end{remark}

\begin{remark}\label{rmk:convention-bridge}
Although \cref{thm:gaussian-fixed} uses the Plancherel normalization (which differs from the probability convention of \cref{def:FT}), the essential content is the same: the Gaussian density and its Fourier transform have the same functional form. In the probability convention, the characteristic function of $\cN(0, \sigma^2)$ is $\varphi(\xi) = e^{-\sigma^2\xi^2/2}$, which is a Gaussian in $\xi$ with ``reciprocal'' variance parameter. The Plancherel normalization is the unique choice that makes the fixed-point identity $\hat{g} = g$ hold \emph{exactly}, without any rescaling.
\end{remark}

\begin{proposition}[{cf.\ \cite[Theorem~VI.5.11]{Parthasarathy}}]\label{prop:gaussian-characterization}
A probability measure $\gamma$ on an LCA group $G$ is Gaussian in the sense of \cref{def:gaussian-lca} if and only if it can be represented as the limit
\[
\gamma = \lim_{n \to \infty} \mu_{n,1} * \mu_{n,2} * \cdots * \mu_{n,k_n}
\]
of convolutions of infinitesimal triangular arrays satisfying appropriate tightness and variance conditions (specifically: the arrays are infinitesimal and uniformly tight; see \cite[Definition~VI.5.1]{Parthasarathy} for precise statements). In particular, Gaussian measures are precisely the possible limits of the CLT on~$G$.
\end{proposition}

%% ============================================================
\section{The heat semigroup and Laplacians on LCA groups}

We now unify the heat equation perspective from \cref{prop:heat-add,prop:heat-mult} in the general LCA framework.

\begin{definition}\label{def:laplacian}
Let $G$ be a connected LCA group with Lie algebra $\mathfrak{g}$. Fix a basis $\{X_1, \ldots, X_d\}$ for $\mathfrak{g}$ and define the \emph{Laplacian} on $G$ by
\[
\Delta_G = \sum_{j=1}^d X_j^2,
\]
where $X_j$ act as left-invariant differential operators on $C^\infty(G)$. The Laplacian depends on the choice of basis (equivalently, on an inner product on $\mathfrak{g}$), but different choices yield the same qualitative theory.
\end{definition}

The characters of $G$ are eigenfunctions of $\Delta_G$. Define the \emph{symbol} $Q : \widehat{G} \to [0, \infty)$ by
\begin{equation}\label{eq:symbol}
\Delta_G \chi = -Q(\chi) \cdot \chi \qquad \text{for all } \chi \in \widehat{G}.
\end{equation}
The function $Q$ is a non-negative definite quadratic form on $\widehat{G}$ in the sense of \cref{def:quadratic}.

\begin{example}\label{ex:laplacian-symbol}
\begin{enumerate}[label=(\alph*)]
\item $G = (\R, +)$: $\Delta_G = \partial_{xx}$, $\chi_\xi(x) = e^{i\xi x}$, $\Delta_G \chi_\xi = -\xi^2 \chi_\xi$, so $Q(\xi) = \xi^2$.
\item $G = (\Rpos, \times)$: $\Delta_G = D^2 = (x\partial_x)^2$, $\chi_s(x) = x^{is}$, $D^2 \chi_s = (is)^2 x^{is} = -s^2 \chi_s$, so $Q(s) = s^2$.
\item $G = (\bT, +)$: $\Delta_G = \partial_{\theta\theta}$, $\chi_n(\theta) = e^{2\pi i n\theta}$, $\Delta_G \chi_n = -(2\pi n)^2 \chi_n$, so $Q(n) = 4\pi^2 n^2$.
\end{enumerate}
\end{example}

\begin{theorem}[Heat semigroup on LCA groups]\label{thm:heat-semigroup}
Let $G$ be a connected LCA group with Laplacian $\Delta_G$ and symbol $Q$ as in \eqref{eq:symbol}. Define the family of Gaussian measures $\{\gamma_t\}_{t > 0}$ by their Fourier transforms:
\[
\widehat{\gamma_t}(\chi) = e^{-tQ(\chi)/2}.
\]
Then:
\begin{enumerate}[label=(\roman*)]
\item $\{\gamma_t\}_{t \geq 0}$ is a convolution semigroup: $\gamma_t * \gamma_s = \gamma_{t+s}$.
\item If $\gamma_t$ has a density $p_t$ with respect to Haar measure, then $p_t$ solves the heat equation
\[
\partial_t p_t = \tfrac{1}{2}\,\Delta_G\, p_t.
\]
\item In Fourier space, the heat equation becomes the family of ODEs (parametrized by $\chi \in \widehat{G}$):
\[
\partial_t \widehat{p_t}(\chi) = -\tfrac{1}{2}\, Q(\chi)\, \widehat{p_t}(\chi),
\]
with solution $\widehat{p_t}(\chi) = e^{-tQ(\chi)/2}$.
\end{enumerate}
\end{theorem}

\begin{proof}
Part (iii) is immediate: applying the Fourier transform to $\partial_t p_t = \frac{1}{2}\Delta_G p_t$ and using $\widehat{\Delta_G p_t}(\chi) = -Q(\chi)\widehat{p_t}(\chi)$ yields the ODE in (iii), whose solution with $\widehat{p_0} = 1$ is $e^{-tQ(\chi)/2}$. Part (i) follows from the multiplicativity $e^{-tQ/2} \cdot e^{-sQ/2} = e^{-(t+s)Q/2}$ and the injectivity of the Fourier transform on probability measures (\cref{thm:levy}). Part (ii) follows from (iii) by inverting the Fourier transform.
\end{proof}

\begin{remark}\label{rmk:heat-clt}
\cref{thm:heat-semigroup} reveals the structural reason the CLT limit is Gaussian: a random walk on $G$, when rescaled, converges to Brownian motion on $G$, whose transition kernel is the heat kernel $p_t$. The Fourier transform diagonalizes $\Delta_G$, reducing the heat PDE to a family of ODEs---one for each character $\chi \in \widehat{G}$---whose solutions are the Gaussian Fourier coefficients $e^{-tQ(\chi)/2}$.
\end{remark}

\begin{example}\label{ex:heat-summary}
Summary of the heat semigroup on the three running examples.
\begin{center}
\renewcommand{\arraystretch}{1.4}
\footnotesize
\begin{tabular}{@{}lcccc@{}}
\textbf{Group $G$} & \textbf{Laplacian} & \textbf{Symbol $Q$} & \textbf{Heat equation} & \textbf{Fund.\ solution} \\
\hline
$(\R, +)$ & $\partial_{xx}$ & $\xi^2$ & $u_t = \tfrac{1}{2}\, u_{xx}$ & Gaussian \\
$(\Rpos, \times)$ & $(x\partial_x)^2$ & $s^2$ & $u_t = \tfrac{1}{2}D^2 u$ & Log-normal \\
$(\bT, +)$ & $\partial_{\theta\theta}$ & $4\pi^2 n^2$ & $u_t = \tfrac{1}{2}\, u_{\theta\theta}$ & Wrapped normal ($\vartheta_3$)
\end{tabular}
\normalsize
\end{center}
\end{example}

%% ============================================================
\section{Ergodic interpretation: ensemble versus time averages}

The CLT admits two complementary interpretations, corresponding to the two sides of the ergodic theorem.

\begin{theorem}[Birkhoff Ergodic Theorem {\cite[Theorem~6.2.1]{Durrett}}]\label{thm:birkhoff}
Let $(\Omega, \mathcal{B}, \mathbb{P}, T)$ be a measure-preserving dynamical system with $T$ ergodic. For any $f \in L^1(\mathbb{P})$,
\[
\frac{1}{n}\sum_{k=0}^{n-1} f(T^k \omega) \longrightarrow \int_\Omega f\, d\mathbb{P} \qquad \text{for $\mathbb{P}$-a.e.\ $\omega$}.
\]
\end{theorem}

\begin{remark}[i.i.d.\ sequences as ergodic systems]\label{rmk:iid-ergodic}
An i.i.d.\ sequence $X_1, X_2, \ldots$ with common distribution $\mu$ can be realized on the product space $\Omega = \R^{\mathbb{N}}$ with the left shift $T(\omega_1, \omega_2, \ldots) = (\omega_2, \omega_3, \ldots)$ and the product measure $\mathbb{P} = \mu^{\otimes \mathbb{N}}$. The Kolmogorov zero-one law implies that $T$ is ergodic. Setting $f(\omega) = \omega_1$, the Birkhoff theorem recovers the strong law of large numbers: $\bar{X}_n \to E[X_1]$ a.s.
\end{remark}

\begin{remark}[Two views of the CLT]\label{rmk:two-views}
The CLT for $\bar{X}_n = \frac{1}{n}\sum_{k=1}^n X_k$ admits two complementary readings:
\begin{enumerate}[label=(\alph*)]
\item \textbf{Ensemble (static, spatial).} Fix $n$ and vary the realization $\omega \in \Omega$. The CLT describes the fluctuations of $\bar{X}_n$ around $\mu$: the distribution of $\sqrt{n}(\bar{X}_n - \mu)$ converges to $\cN(0, \sigma^2)$ as $n \to \infty$.

\item \textbf{Temporal (dynamic).} Fix a typical $\omega$ and let $n \to \infty$. The LLN says $\bar{X}_n(\omega) \to \mu$ a.s., describing convergence along a single trajectory.
\end{enumerate}
The ensemble CLT describes the \emph{shape} of the fluctuations (Gaussian), while the ergodic theorem guarantees that time averages equal space averages. These are dual aspects of the same structure: the heat semigroup provides the universal transition kernel, and the Fourier transform diagonalizes it whether we sum over space or over time.
\end{remark}

\begin{remark}[Statistical mechanics]\label{rmk:statmech}
In the language of statistical mechanics, the ensemble average corresponds to the Gibbs measure (a spatial average over configurations), while the time average corresponds to trajectory averages of a single system. The equality of the two is the content of ergodicity. The Gaussian arising from the CLT is the universal fluctuation law in both pictures. On the multiplicative group, the log-normal distribution governs the analogous fluctuations---for instance, the growth rates in multiplicative processes or the fluctuations of products of random matrices' singular values.
\end{remark}

%% ============================================================
\section{Hardy spaces and causal projections}\label{sec:hardy}

The Fourier--Mellin duality of \cref{thm:mellin-fourier} acts on all of $L^2(G)$. But $L^2(G)$ carries a finer structure: an orthogonal decomposition into ``causal'' (forward-propagating) and ``anti-causal'' halves, governed by the \emph{Hardy space}. We show that the $\log$ isomorphism preserves this decomposition, relate the CLT to prediction theory via the Wold decomposition of the time-shift isometry, and identify the Toeplitz algebra as the natural operator-algebraic framework.

\begin{definition}[Hardy spaces on LCA groups]\label{def:hardy-space}
For each of the three running examples of \cref{ex:groups}, define the \emph{Hardy space} $\Htwo(G) \subset L^2(G)$ as follows.
\begin{enumerate}[label=(\alph*)]
\item \textbf{$G = (\R, +)$.} The Hardy space of the real line is
\[
\Htwo(\R) = \bigl\{f \in L^2(\R) : \widehat{f}(\xi) = 0 \text{ for a.e.\ } \xi < 0\bigr\}.
\]
Equivalently, $f \in \Htwo(\R)$ if and only if $f$ extends to a holomorphic function in the upper half-plane $\mathbb{C}_+ = \{z : \operatorname{Im} z > 0\}$ with
$\sup_{y > 0} \int_{-\infty}^{\infty} |f(x + iy)|^2\, dx < \infty$
\cite[Chapter~II]{Garnett}. These are the \emph{causal} signals: their spectral content is one-sided.

\item \textbf{$G = (\Rpos, \times)$.} Define the \emph{multiplicative Hardy space} as the preimage of $\Htwo(\R)$ under the $\log$ isomorphism:
\[
\Htwo(\Rpos) = \bigl\{f \in L^2(\Rpos, dx/x) : f \circ \exp \in \Htwo(\R)\bigr\}.
\]
Equivalently, $f \in \Htwo(\Rpos)$ if and only if its Mellin transform $\cM[f](s)$ extends analytically to $\operatorname{Im}(s) > 0$ with uniformly bounded $L^2$ norms on horizontal lines \cite[Chapter~4]{Hoffman}. This is the \emph{scale-causal} subspace.

\item \textbf{$G = (\bT, +)$.} The classical Hardy space on the circle is
\[
\Htwo(\bT) = \bigl\{f \in L^2(\bT) : \widehat{f}(n) = 0 \text{ for all } n < 0\bigr\}.
\]
Equivalently, $f$ extends to a holomorphic function on the open unit disk $\mathbb{D}$ with $\sup_{0 < r < 1} \int_0^1 |f(re^{2\pi i\theta})|^2\, d\theta < \infty$.
\end{enumerate}

In each case, the ``positive frequencies'' are determined by a choice of ordering on $\widehat{G}$: the standard ordering on $\widehat{G} \cong \R$ (for $G = (\R,+)$ and $G = (\Rpos,\times)$) and the standard ordering on $\widehat{G} = \Z$ (for $G = (\bT,+)$). This ordering encodes a choice of \emph{time direction}---the causality arrow.
\end{definition}

\begin{proposition}[Hardy decomposition]\label{prop:hardy-decomp}
For each $G \in \{(\R, +),\, (\Rpos, \times),\, (\bT, +)\}$, the Hilbert space $L^2(G)$ admits the orthogonal decomposition
\begin{equation}\label{eq:hardy-decomp}
L^2(G) = \Htwo(G) \oplus \Htwo(G)^\perp,
\end{equation}
where $\Htwo(G)^\perp = \{f \in L^2(G) : \widehat{f}(\chi) = 0 \text{ for } \chi \geq 0\}$. The orthogonal projection $P_{\Htwo} : L^2(G) \to \Htwo(G)$ is given in the frequency domain by
\[
\widehat{P_{\Htwo} f}(\chi) = \mathbf{1}_{\chi \geq 0}\, \widehat{f}(\chi).
\]
\end{proposition}

\begin{proof}
Write $f = f^+ + f^-$ with $\widehat{f^+} = \mathbf{1}_{\chi \geq 0}\,\widehat{f}$ and $\widehat{f^-} = \mathbf{1}_{\chi < 0}\,\widehat{f}$. Then $f^+ \in \Htwo(G)$, $f^- \in \Htwo(G)^\perp$, and by Parseval's theorem,
\[
\langle f^+, f^- \rangle = \int_{\widehat{G}} \widehat{f^+}(\chi)\, \overline{\widehat{f^-}(\chi)}\, d\chi = 0,
\]
since the supports of $\widehat{f^+}$ and $\widehat{f^-}$ are disjoint.
\end{proof}

\begin{theorem}[Fourier--Mellin duality preserves Hardy structure]\label{thm:hardy-duality}
The isometric isomorphism $\log_* : L^2(\Rpos, dx/x) \to L^2(\R, dy)$ of \cref{prop:log-iso} restricts to an isometric isomorphism
\[
\log_* : \Htwo(\Rpos) \xrightarrow{\;\sim\;} \Htwo(\R),
\]
and conjugates the Mellin-Hardy projection to the Fourier-Hardy projection:
\[
\log_* \circ P_{\Htwo(\Rpos)} = P_{\Htwo(\R)} \circ \log_*.
\]
That is, the following diagram commutes:
\begin{equation}\label{eq:cd-hardy}
\begin{tikzcd}
L^2(\Rpos, dx/x) \arrow[r, "P_{\Htwo(\Rpos)}"] \arrow[d, "\log_*"'] & \Htwo(\Rpos) \arrow[d, "\log_*"] \\
L^2(\R, dy) \arrow[r, "P_{\Htwo(\R)}"'] & \Htwo(\R)
\end{tikzcd}
\end{equation}
This lifts \cref{thm:mellin-fourier} from the full $L^2$ level to the Hardy subspace level.
\end{theorem}

\begin{proof}
By \cref{thm:mellin-fourier}, $\cM[f](s) = \cF[\log_* f](s)$ for every $s \in \R$. In particular, $\cM[f](s) = 0$ for $s < 0$ if and only if $\cF[\log_* f](s) = 0$ for $s < 0$, i.e., $f \in \Htwo(\Rpos)$ if and only if $\log_* f \in \Htwo(\R)$. Since $\log_*$ is an isometry (it preserves Haar measure: $dx/x \mapsto dy$), the restriction is an isometric isomorphism. The projection intertwining follows: for any $f \in L^2(\Rpos, dx/x)$,
\[
\log_*(P_{\Htwo(\Rpos)} f) = P_{\Htwo(\R)}(\log_* f),
\]
because both sides have the same Fourier transform $\mathbf{1}_{s \geq 0}\, \cF[\log_* f](s)$.
\end{proof}

\begin{definition}[Unilateral shift]\label{def:shift}
On $\Htwo(\bT)$, define the \emph{unilateral shift} $S : \Htwo(\bT) \to \Htwo(\bT)$ by
\begin{equation}\label{eq:shift}
(Sf)(z) = z\, f(z), \qquad \text{equivalently,} \qquad \widehat{(Sf)}(n) = \widehat{f}(n-1) \text{ for } n \geq 0.
\end{equation}
Its adjoint satisfies $\widehat{(S^*f)}(n) = \widehat{f}(n+1)$, whence:
\begin{itemize}
\item $S^*S = I$ \quad (S is an isometry),
\item $SS^* = I - \langle \,\cdot\,, \mathbf{1}\rangle\, \mathbf{1} \neq I$ \quad (S is not unitary),
\item $\ker(S^*) = \operatorname{span}\{\mathbf{1}\}$ \quad (the constant functions).
\end{itemize}
On $\Htwo(\R)$, the analogue is the half-line translation $T_a$ ($a > 0$) defined by $\widehat{T_a f}(\xi) = e^{ia\xi}\widehat{f}(\xi)$ for $\xi \geq 0$. On $\Htwo(\Rpos)$, the Mellin-side analogue is the dilation $S_a$ defined by $\widehat{S_a g}(s) = a^{is}\widehat{g}(s)$ for $s \geq 0$. By \cref{thm:hardy-duality}, $\log_*$ conjugates $S_a$ to $T_a$.
\end{definition}

\begin{theorem}[Wold decomposition {\cite[Chapter~I]{SzNF}}]\label{thm:wold}
Let $V$ be an isometry on a Hilbert space $\mathcal{H}$. Then
\begin{equation}\label{eq:wold}
\mathcal{H} = \left(\bigoplus_{n=0}^{\infty} V^n(\ker V^*)\right) \oplus \mathcal{H}_u, \qquad \mathcal{H}_u = \bigcap_{n \geq 0} V^n \mathcal{H},
\end{equation}
where $\mathcal{H}_u$ is the maximal reducing subspace on which $V$ is unitary. The first summand is the \emph{purely non-deterministic} part; the second is the \emph{deterministic} part. For the unilateral shift $S$ on $\Htwo(\bT)$, one has $\mathcal{H}_u = \{0\}$: the shift is purely non-deterministic.
\end{theorem}

\begin{proof}[Proof sketch]
Since $V$ is an isometry, $V^n(\ker V^*)$ and $V^m(\ker V^*)$ are orthogonal for $n \neq m$: if $n > m$, then $V^n(\ker V^*) \subset V^{n}\mathcal{H}$ while $V^m(\ker V^*) \perp V^{m+1}\mathcal{H} \supset V^n\mathcal{H}$. Define $\mathcal{H}_u = \bigcap_{n \geq 0} V^n\mathcal{H}$; this subspace is $V$-invariant, and $V|_{\mathcal{H}_u}$ is surjective (hence unitary) since $V(\mathcal{H}_u) = V\bigl(\bigcap V^n\mathcal{H}\bigr) = \bigcap V^{n+1}\mathcal{H} = \mathcal{H}_u$. For $S$ on $\Htwo(\bT)$: $S^n\Htwo = z^n\Htwo$ and $\bigcap_{n \geq 0} z^n\Htwo = \{0\}$ by the $F.\ $and $M.\ $Riesz theorem.
\end{proof}

Each layer $V^n(\ker V^*)$ represents information that becomes causally inaccessible after $n$ steps of forward evolution. The decomposition separates \emph{unpredictable innovations} (the direct sum) from the \emph{perfectly predictable} component ($\mathcal{H}_u$).

\begin{proposition}[CLT and prediction theory]\label{prop:clt-prediction}
Let $\{X_n\}_{n \in \Z}$ be a wide-sense stationary process with spectral density $f \in L^1(\bT)$, $f \geq 0$. Let $V$ act on $\overline{\operatorname{span}}\{X_n : n \in \Z\} \subset L^2(\Omega)$ as the time shift $VX_n = X_{n+1}$.
\begin{enumerate}[label=(\roman*)]
\item \textbf{(Szeg\H{o}--Kolmogorov.)} The process is purely non-deterministic ($\mathcal{H}_u = \{0\}$ in \cref{thm:wold}) if and only if the \emph{Szeg\H{o} condition} holds:
\[
\int_0^1 \log f(\theta)\, d\theta > -\infty.
\]

\item \textbf{(Innovation and outer function.)} When purely non-deterministic, $f$ admits an outer function factorization $f(\theta) = |h(e^{2\pi i\theta})|^2$ with $h \in \Htwo(\bT)$ outer, and the one-step prediction error is
\[
\sigma^2 = \exp\left(\int_0^1 \log f(\theta)\, d\theta\right).
\]

\item \textbf{(CLT link.)} For i.i.d.\ sequences, $f \equiv \sigma^2$ is constant, the Szeg\H{o} condition is trivially satisfied, $\mathcal{H}_u = \{0\}$, and the CLT variance is $\sigma^2$. More generally, the CLT for partial sums of stationary processes holds under Gordin's condition \cite[Theorem~35.5]{Billingsley}, which strengthens pure non-determinism.
\end{enumerate}
By \cref{thm:hardy-duality}, the entire prediction-theoretic structure transports to $(\Rpos, \times)$: the multiplicative CLT for products of positive random variables has the same Wold decomposition and Szeg\H{o} condition, with the Mellin spectral density replacing the Fourier spectral density.
\end{proposition}

\begin{remark}[Toeplitz algebra and Fredholm index]
For $\varphi \in C(\bT)$, define the \emph{Toeplitz operator} $T_\varphi : \Htwo(\bT) \to \Htwo(\bT)$ by $T_\varphi f = P_{\Htwo}(\varphi f)$. The \emph{Toeplitz algebra} $\cT = C^*(T_\varphi : \varphi \in C(\bT))$ fits into the short exact sequence \cite[Chapter~7]{Douglas}
\[
0 \to \cK \to \cT \to C(\bT) \to 0,
\]
where $\cK$ denotes the compact operators on $\Htwo(\bT)$. When $\varphi$ is non-vanishing, $T_\varphi$ is Fredholm with index \cite[Chapter~7]{Nikolski}
\[
\ind(T_\varphi) = -\operatorname{wind}(\varphi, 0),
\]
where $\operatorname{wind}(\varphi, 0)$ is the winding number of $\varphi$ around the origin. The Fourier--Mellin duality of \cref{thm:hardy-duality} induces a corresponding isomorphism of \emph{Mellin--Toeplitz algebras} on $\Htwo(\Rpos)$: defining $T_\varphi^{\cM} g = P_{\Htwo(\Rpos)}(\varphi \cdot g)$, the $\log$ isomorphism conjugates $T_\varphi^{\cM}$ to $T_{\varphi \circ \exp}$, preserving the Fredholm index since $\log$ (being a homeomorphism) preserves winding numbers.
\end{remark}

\begin{definition}[Toeplitz determinant and geometric mean]\label{def:toeplitz-det}
For $f \in L^1(\bT)$ with $f > 0$ a.e., define:
\begin{enumerate}[label=(\alph*)]
\item The $n \times n$ \emph{Toeplitz matrix} $T_n(f) = (\widehat{f}_{j-k})_{0 \leq j,k \leq n-1}$, i.e., the compression of the multiplication operator $M_f$ to $\operatorname{span}\{1, z, \ldots, z^{n-1}\} \subset \Htwo(\bT)$.
\item The \emph{Toeplitz determinant} $D_n(f) = \det T_n(f)$.
\item The \emph{geometric mean} $G(f) = \exp\bigl(\int_0^1 \log f(\theta)\, d\theta\bigr)$, well-defined and positive provided $\log f \in L^1(\bT)$.
\end{enumerate}
Note that $D_n(f)$ is a \emph{multiplicative} quantity---the product of eigenvalues of $T_n(f)$, interpreted as the volume of an $n$-dimensional parallelepiped in the ``multiplicative theater.'' The quantity $G(f)$ is the exponential of an \emph{additive} average: $\int \log f$ lives in the ``additive theater.'' The Szeg\H{o} theorems below bridge these two theaters with a rigid, non-trivializing comparison.
\end{definition}

\begin{theorem}[Szeg\H{o} Limit Theorems {\cite[Chapter~5]{BottcherSilbermann}}]\label{thm:szego}
Let $f \in L^1(\bT)$ with $f > 0$ a.e.\ and $\log f \in L^1(\bT)$. Write $c_k = \int_0^1 (\log f(\theta))\, e^{-2\pi i k\theta}\, d\theta$ for the Fourier coefficients of $\log f$.
\begin{enumerate}[label=(\roman*)]
\item \textbf{(First Szeg\H{o} Limit Theorem, 1915.)} The geometric mean of the eigenvalues of $T_n(f)$ converges to the exponential of the arithmetic mean of $\log f$:
\begin{equation}\label{eq:szego-first}
\lim_{n \to \infty} D_n(f)^{1/n} = G(f).
\end{equation}

\item \textbf{(Strong Szeg\H{o} Limit Theorem, 1952.)} If additionally $\sum_{k \in \Z} |k|\, |c_k|^2 < \infty$, then
\begin{equation}\label{eq:szego-strong}
\lim_{n \to \infty} \frac{D_n(f)}{G(f)^n} = E(f),
\end{equation}
where the \emph{Szeg\H{o} constant} is
\begin{equation}\label{eq:szego-constant}
E(f) = \exp\!\left(\sum_{k=1}^{\infty} k\, |c_k|^2\right) \in (0, \infty).
\end{equation}
\end{enumerate}
\end{theorem}

\begin{proof}[Proof sketch]
For~(i): the eigenvalue distribution of $T_n(f)$ converges weakly to $f$ (Szeg\H{o} distribution theorem), so $(1/n)\sum_{j=1}^n \log \lambda_j^{(n)} \to \int_0^1 \log f(\theta)\, d\theta$, giving $D_n(f)^{1/n} \to G(f)$.

For~(ii): write $\log f = c_0 + h + \bar{h}$ where $h = \sum_{k=1}^{\infty} c_k e^{2\pi i k \theta} \in \Htwo(\bT)$ is the \emph{causal part} of $\log f$ --- its Hardy projection (\cref{prop:hardy-decomp}). Then $f = G(f) \cdot |e^h|^2$, and the factorization $f = |G(f)^{1/2}\, e^h|^2$ is precisely the outer function factorization of \cref{prop:clt-prediction}(ii). The Toeplitz determinant factors as
\[
D_n(f) = G(f)^n \cdot \det\bigl(T_n(e^h)\, T_n(e^{\bar{h}})\bigr),
\]
and the Borodin--Okounkov--Case--Geronimo identity \cite[Chapter~5]{BottcherSilbermann} gives $\det(T_n(e^h)\, T_n(e^{\bar{h}})) \to \exp(\sum_{k=1}^{\infty} k\, |c_k|^2)$. Alternatively, see \cite[Chapter~2]{Simon}.
\end{proof}

The factorization $\log f = c_0 + h + \bar{h}$ with $h \in \Htwo(\bT)$ is the Hardy decomposition of $\log f$. The causal part $h$ generates the outer function; the constant $c_0 = \int \log f$ generates $G(f) = e^{c_0}$. The Szeg\H{o} constant $E(f)$ therefore measures the ``interaction energy'' between the causal and anti-causal components of $\log f$.

\begin{corollary}[Non-trivial cross-theater volume comparison]\label{cor:nontrivial-volume}
Under the hypotheses of \cref{thm:szego}(ii),
\[
D_n(f) = G(f)^n \cdot E(f) \cdot (1 + o(1)), \qquad 0 < E(f) < \infty.
\]
In particular:
\begin{enumerate}[label=(\roman*)]
\item The \textbf{multiplicative} volume $D_n(f)$ (determinant, living in the multiplicative theater) is asymptotically determined by the \textbf{additive} quantity $\int \log f$ (living in the additive theater), up to a finite, nonzero, computable correction factor $E(f)$.
\item The comparison does \textbf{not} trivialize: neither $E(f) = 0$ nor $E(f) = \infty$.
\item The analogue for non-compact groups proceeds in two steps. First, on $(\R, +)$, the Toeplitz matrices $T_n(f)$ are replaced by truncated \emph{Wiener--Hopf operators} $W_R(a)$ on $L^2(0, R)$, and the Szeg\H{o} asymptotics carry over with $n \to \infty$ replaced by $R \to \infty$ \cite[Chapter~10]{BottcherSilbermann}. Second, the Hardy duality (\cref{thm:hardy-duality}) transports this to $(\Rpos, \times)$: the Mellin-side Wiener--Hopf operator on $\Htwo(\Rpos)$ satisfies the same asymptotics with the same Szeg\H{o} constant, since $\log_*$ is an isometry preserving all spectral data.
\item For i.i.d.\ sequences ($f \equiv \sigma^2$ constant), $c_k = 0$ for all $k \neq 0$, hence $E(f) = 1$ exactly. The Gaussian is the unique spectral density for which the cross-theater comparison is exact---no correction.
\end{enumerate}
\end{corollary}

\begin{remark}[The Szeg\H{o} condition as universal constraint]
The condition $\int_{\bT} \log f(\theta)\, d\theta > -\infty$ has now appeared in three apparently independent roles:
\begin{enumerate}
\item \textbf{Prediction theory} (\cref{prop:clt-prediction}(i)): necessary and sufficient for pure non-determinism ($\mathcal{H}_u = \{0\}$).
\item \textbf{Central limit theorem} (\cref{prop:clt-prediction}(iii)): necessary for the CLT to hold via Gordin's martingale method for stationary processes with standard $\sqrt{n}$ normalization.
\item \textbf{Non-trivial volume comparison} (\cref{thm:szego,cor:nontrivial-volume}): necessary for $G(f) \in (0, \infty)$, i.e., for the additive-to-multiplicative bridge to be well-defined.
\end{enumerate}
This triple coincidence is structural. The Gaussian is simultaneously the Fourier eigenfunction (\cref{thm:gaussian-fixed}), the heat kernel (\cref{thm:heat-semigroup}), the outer function seed ($c_k = 0$ for $k \neq 0$), and the unique spectral density for which $E(f) = 1$ exactly---\emph{perfect additive--multiplicative alignment}.

The CLT identifies the Gaussian as the universal \emph{attractor}: short-memory stationary processes converging to the Gaussian limit generically satisfy the Szeg\H{o} condition $\int \log f > -\infty$ (ensuring $G(f) \in (0,\infty)$ and the First Szeg\H{o} Theorem). The Strong Szeg\H{o} regularity $\sum k|c_k|^2 < \infty$ is an \emph{additional} condition---not implied by the CLT per se, but naturally satisfied in its vicinity. Precisely: the class $\{f : \sum k|c_k|^2 < \infty\}$ is an open neighborhood of the constant density $f \equiv \sigma^2$ (the Gaussian fixed point) in the topology of $L^1(\bT)$, and any spectral density with $\log f \in H^{1/2}(\bT)$---the Sobolev space of half-order regularity---lies in this class. Short-memory processes satisfying the CLT (summable correlations, continuous spectral density) generically have $\log f$ of this regularity. The Gaussian fixed point thus sits at the center of the \emph{basin of non-triviality}: the region in spectral density space where cross-theater volume comparisons yield finite, nonzero bounds.
\end{remark}

\begin{remark}[Non-trivial versus trivial multiverses]
The Szeg\H{o} Limit Theorems provide a precise framework for comparing quantities across different algebraic universes---the additive theater $(\R, +)$ and the multiplicative theater $(\Rpos, \times)$---while maintaining rigidity. We say that a multiverse of parallel algebraic theaters is \emph{non-trivial} if:
\begin{enumerate}[label=(\alph*)]
\item the theaters are connected by \emph{isometric} operators (the Fourier--Mellin duality, \cref{thm:hardy-duality});
\item the per-step information loss is \emph{bounded} ($\dim \ker S^* = 1$, \cref{thm:wold});
\item cross-theater volume comparisons yield \emph{finite, nonzero} correction factors ($E(f) \in (0, \infty)$, \cref{thm:szego}(ii)).
\end{enumerate}
A multiverse \emph{trivializes} when any condition fails: (a${}'$)~the inter-theater link introduces unbounded indeterminacy, destroying isometry; (b${}'$)~the information loss per step is uncontrolled; (c${}'$)~the volume comparison yields $E(f) = \infty$ or $E(f) = 0$, reducing the cross-theater inequality to $\infty \geq \text{const}$---a tautology with no constraining power.

We offer the following structural analogy, without claiming a formal correspondence between the two frameworks. In the inter-universal Teichm\"uller theory of Mochizuki \cite{MochizukiIII}, an analogous comparison of log-volumes across arithmetically linked Hodge theaters is attempted in Corollary~3.12. As analyzed by Scholze and Stix \cite{ScholzeStix}, the indeterminacies introduced to align the theaters---in particular, the unbounded distortion of the log-shell permitted by Indeterminacy~3---play a role structurally parallel to allowing the spectral density $f$ to vary without the constraint $\sum k |c_k|^2 < \infty$: in both cases, the ``inter-universe comparison'' degenerates because the correction factor ($E(f)$ in our setting, the log-shell distortion in IUT) is unbounded. The resulting inequality trivializes to $\infty \geq \mathrm{const}$.

The difference between a trivial and non-trivial multiverse reduces to a single analytic condition: the finiteness of $E(f)$. This condition is equivalent to $\log f \in H^{1/2}(\bT)$---a regularity generically satisfied by short-memory processes in the CLT regime, but logically independent of the CLT itself. The relationship is mediated by the fixed point: the CLT identifies the Gaussian ($f \equiv \sigma^2$, $E(f) = 1$) as the universal attractor, and the Strong Szeg\H{o} class forms an open basin around it. The central limit theorem, absent from the IUT framework, provides both the attractor and the mechanism that draws spectral densities toward it; the Strong Szeg\H{o} condition delineates the basin within which this attraction yields non-trivial cross-theater bounds. In the language of this paper: Mochizuki constructed a multiverse without a fixed point. The Fourier--Mellin duality (\cref{thm:mellin-fourier}), the Hardy isometry (\cref{thm:hardy-duality}), and the CLT convergence (\cref{prop:clt-prediction}) together ensure that our multiverse \emph{has} a fixed point---the Gaussian---and the Strong Szeg\H{o} regularity, satisfied at and around this fixed point, forces $E(f) < \infty$.
\end{remark}

\begin{remark}[Connection to viability theory]
The Hardy decomposition $L^2 = \Htwo \oplus (\Htwo)^\perp$ has a dynamical-systems interpretation connecting to viability theory \cite{Aubin}. The Hardy space $\Htwo$ consists of modes whose spectral support is ``forward''---the \emph{causal} or \emph{viable} modes:
\begin{itemize}
\item The viability kernel $\operatorname{Viab}(K)$ of a constraint set $K$ under a differential inclusion $\dot{x} \in F(x)$ \cite[Chapter~4]{Aubin} consists of the initial states from which forward evolution remains in $K$. In the spectral setting, $\Htwo$ plays this role: it consists of functions whose frequency content permits causal (forward-time) evolution.
\item The spectral gap of $\Delta_G$---equivalently, the CLT convergence rate (cf.\ \cref{thm:heat-semigroup})---governs the rate at which viable trajectories approach the Gaussian equilibrium.
\item The Wold decomposition's per-step information loss $\ker(V^*)$ captures the irreversibility of causal projection: information about anti-causal modes is lost under forward propagation, analogous to the loss of information across a causal horizon.
\item The Toeplitz exact sequence $0 \to \cK \to \cT \to C(\bT) \to 0$ encodes the passage from operator-algebraic to commutative descriptions---tracing out the ``anti-causal'' degrees of freedom.
\end{itemize}
\end{remark}

\begin{example}[Summary: Hardy spaces on the three groups]\label{ex:hardy-summary}
\[
\renewcommand{\arraystretch}{1.4}
\begin{array}{lllll}
\hline
\textbf{Group } G & \Htwo(G) & \textbf{Shift } S & \ker(S^*) & \textbf{Toeplitz symbol} \\
\hline
(\R, +) & \widehat{f}(\xi) = 0,\; \xi < 0 & T_a:\; e^{ia\xi}\widehat{f}(\xi) & \text{---} & L^\infty(\R) \cap C_0 \\
(\Rpos, \times) & \cM[f](s) = 0,\; s < 0 & S_a:\; a^{is}\widehat{g}(s) & \text{---} & L^\infty(\R) \cap C_0 \\
(\bT, +) & \widehat{f}(n) = 0,\; n < 0 & z\text{-multiplication} & \operatorname{span}\{\mathbf{1}\} & C(\bT) \\
\hline
\end{array}
\]
The second row is the image of the first under $\exp_*$ (equivalently, the first is the image of the second under $\log_*$), reflecting \cref{thm:hardy-duality}. Here $C_0$ denotes continuous functions vanishing at infinity. On the non-compact groups $(\R, +)$ and $(\Rpos, \times)$, the shift is a one-parameter semigroup $\{T_a\}_{a > 0}$ (resp.\ $\{S_a\}_{a > 1}$) of isometries rather than a single isometry, so the Wold decomposition takes the continuous form (Lax--Phillips scattering theory) and $\ker(S^*)$ is replaced by the infinitesimal generator's defect space---marked ``---'' since it does not reduce to a finite-dimensional subspace.
\end{example}

%% ============================================================
\section{Concluding remarks}

\subsection*{The exponential map as the general bridge}

The Fourier--Mellin duality of Sections~3--5 rests on the group isomorphism $\exp : (\R, +) \to (\Rpos, \times)$, which bridges the ``additive'' world (Lie algebra) and the ``multiplicative'' world (Lie group). This bridge generalizes: every Lie group $G$ with Lie algebra $\mathfrak{g} = T_e G$ admits an \emph{exponential map}
\[
\exp : \mathfrak{g} \to G, \qquad \exp(X) = \gamma_X(1),
\]
where $\gamma_X : \R \to G$ is the unique one-parameter subgroup satisfying $\gamma_X'(0) = X$ \cite[Chapter~3]{Hall}. A crucial distinction separates the abelian and non-abelian settings:
\begin{itemize}
\item \textbf{Abelian case.} The map $\exp : (\R, +) \to (\Rpos, \times)$ is a \emph{global} isomorphism. The Fourier--Mellin duality is exact at all scales.
\item \textbf{Non-abelian case.} The map $\exp : \mathfrak{g} \to G$ is in general only a \emph{local} diffeomorphism near the origin. Globally, it may fail to be injective (due to periodicity, as in $SO(3)$), surjective (as in $SL(2, \R)$), or both.
\end{itemize}
The CLT, however, concerns \emph{small fluctuations} near the identity---precisely the regime where $\exp$ is a diffeomorphism. The general picture is:

\begin{center}
\begin{tikzcd}[column sep=large]
\text{Gaussian on } \mathfrak{g} \arrow[r, "\exp_*"] & \text{CLT limit on } G \text{ (near $e$)}
\end{tikzcd}
\end{center}

\noindent This generalizes \cref{cor:clt-image}: just as $\operatorname{LogNormal} = \exp_*\cN$ on $\Rpos$, the CLT limit on any Lie group is locally the pushforward of a Gaussian on $\mathfrak{g}$ under $\exp$.

\subsection*{Non-abelian groups: $GL(n)$ and the Cartan decomposition}

On non-abelian groups such as $GL(n, \R)$, the Pontryagin dual is replaced by the unitary dual (irreducible unitary representations), and the Fourier transform becomes operator-valued. For products of i.i.d.\ random matrices $M_n = A_1 A_2 \cdots A_n$, the CLT takes a decomposed form via the Cartan ($KAK$) decomposition $GL(n) \cong K \cdot A \cdot K$, where $K = O(n)$ and $A$ consists of positive diagonal matrices:
\begin{enumerate}[label=(\roman*)]
\item The ``angular'' part (in $K = O(n)$) converges to Haar measure on $O(n)$---a mixing phenomenon.
\item The ``radial'' part---the Lyapunov exponents $\lambda_1 \geq \cdots \geq \lambda_n$ parametrizing $A$---satisfies a standard Gaussian CLT on the Cartan subalgebra $\mathfrak{a} \cong \R^n$ \cite{BQ}.
\end{enumerate}
After centering, the fluctuations lie in the hyperplane $\sum_i (\lambda_i - \bar{\lambda}) = 0$, which is precisely the Lie algebra $\mathfrak{sl}(n) = \{X \in \mathfrak{gl}(n) : \tr(X) = 0\}$. Thus the centered CLT limit on $GL(n)$ is a Gaussian measure on $\mathfrak{sl}(n)$---not the group $SL(n)$ itself, but its tangent space at the identity. The passage from $\mathfrak{sl}(n)$ to $SL(n)$ is exactly the exponential map $\exp : \mathfrak{sl}(n) \to SL(n)$, extending the abelian bridge $\exp : \R \to \Rpos$ to the matrix setting.

\subsection*{Compact groups: closed-form exponential maps}

On compact groups, the exponential map admits elegant closed forms that make the passage $\mathfrak{g} \to G$ explicit.

\medskip\noindent\textbf{$\boldsymbol{SO(3)}$ and Rodrigues' formula.}
The Lie algebra $\mathfrak{so}(3)$ consists of $3 \times 3$ skew-symmetric matrices. Via the \emph{hat map} $\widehat{(\cdot)} : \R^3 \to \mathfrak{so}(3)$ defined by $\hat{\mathbf{n}}\mathbf{v} = \mathbf{n} \times \mathbf{v}$, every element of $\mathfrak{so}(3)$ can be written as $\theta\hat{\mathbf{n}}$ for a unit vector $\mathbf{n} \in S^2$ and an angle $\theta \in \R$.

The key algebraic fact is that $\hat{\mathbf{n}}$ (for $\|\mathbf{n}\| = 1$) has eigenvalues $0, +i, -i$, so by the Cayley--Hamilton theorem, it satisfies its characteristic polynomial:
\begin{equation}\label{eq:cayley-hamilton-so3}
\hat{\mathbf{n}}^3 + \hat{\mathbf{n}} = 0, \qquad\text{i.e.,}\quad \hat{\mathbf{n}}^3 = -\hat{\mathbf{n}}.
\end{equation}
This periodicity causes the matrix exponential series to collapse. Separating even and odd powers:
\begin{align*}
\exp(\theta\hat{\mathbf{n}})
&= \sum_{k=0}^{\infty} \frac{(\theta\hat{\mathbf{n}})^k}{k!} \\
&= I + \underbrace{\left(\theta - \frac{\theta^3}{3!} + \frac{\theta^5}{5!} - \cdots\right)}_{\displaystyle\sin\theta}\hat{\mathbf{n}} \;+\; \underbrace{\left(\frac{\theta^2}{2!} - \frac{\theta^4}{4!} + \frac{\theta^6}{6!} - \cdots\right)}_{\displaystyle 1 - \cos\theta}\hat{\mathbf{n}}^2.
\end{align*}
This yields the \emph{Rodrigues rotation formula}:
\begin{equation}\label{eq:rodrigues}
\exp(\theta\hat{\mathbf{n}}) = I + \sin\theta\;\hat{\mathbf{n}} + (1 - \cos\theta)\;\hat{\mathbf{n}}^2.
\end{equation}
The infinite series reduces to three terms because the minimal polynomial of $\hat{\mathbf{n}}$ has degree~3; the odd-power terms resum as $\sin\theta$ and the even-power terms as $1 - \cos\theta$.

\medskip\noindent\textbf{$\boldsymbol{SU(2)}$, quaternions, and the double cover.}
The group $SU(2)$ of $2 \times 2$ unitary matrices with determinant $1$ is diffeomorphic to $S^3 \subset \R^4$ via identification with unit quaternions. Its Lie algebra $\mathfrak{su}(2)$ is identified with the pure imaginary quaternions $\operatorname{Im}(\mathbb{H}) \cong \R^3$.

For a unit vector $\mathbf{n} \in S^2 \subset \operatorname{Im}(\mathbb{H})$, the relation $\mathbf{n}^2 = -1$ (the defining property of imaginary units in $\mathbb{H}$) yields the exponential map via direct application of Euler's formula:
\begin{equation}\label{eq:quat-exp}
\exp\!\left(\tfrac{\theta}{2}\,\mathbf{n}\right) = \cos\tfrac{\theta}{2} + \sin\tfrac{\theta}{2}\;\mathbf{n} \;\in\; S^3 \cong SU(2).
\end{equation}
This is the quaternionic generalization of $e^{i\alpha} = \cos\alpha + i\sin\alpha$: the proof is identical, with $\mathbf{n}^2 = -1$ replacing $i^2 = -1$.

The factor $\frac{1}{2}$ reflects the double cover $\pi : SU(2) \to SO(3)$, defined by $\pi(q) : \mathbf{v} \mapsto q\mathbf{v}\bar{q}$ for $\mathbf{v} \in \operatorname{Im}(\mathbb{H}) \cong \R^3$. Since $q$ and $-q$ act identically by conjugation, $\ker(\pi) = \{\pm 1\}$. Composing \eqref{eq:quat-exp} with $\pi$ recovers the Rodrigues formula \eqref{eq:rodrigues}: the quaternion $\cos\frac{\theta}{2} + \sin\frac{\theta}{2}\,\mathbf{n}$ maps to the rotation matrix $I + \sin\theta\,\hat{\mathbf{n}} + (1-\cos\theta)\hat{\mathbf{n}}^2$, as can be verified by expanding $q\mathbf{v}\bar{q}$ and using the double-angle identities $\sin\theta = 2\sin\frac{\theta}{2}\cos\frac{\theta}{2}$ and $1 - \cos\theta = 2\sin^2\frac{\theta}{2}$.

\medskip\noindent\textbf{CLT on compact groups.}
Since $SO(3)$ and $SU(2)$ are compact, repeated convolutions converge to Haar measure (the uniform distribution on $G$) rather than to a Gaussian: a compact group has finite volume, leaving no room for the unbounded fluctuations that a Gaussian requires. However, for \emph{small times}---when the random increments are small perturbations of the identity---the heat kernel $p_t$ on $G$ is well-approximated by the Gaussian on $\mathfrak{g}$ pushed forward by $\exp$:
\[
p_t(g) \;\approx\; \frac{1}{(2\pi t)^{d/2}}\, \exp\!\left(-\frac{\|\exp^{-1}(g)\|^2}{2t}\right) \qquad \text{as } t \to 0^+,\; g \to e,
\]
where $d = \dim G$ and $\|\cdot\|$ is the norm on $\mathfrak{g}$ induced by the Riemannian metric. As $t \to \infty$, the heat kernel converges to the constant $1/\operatorname{vol}(G)$ (Haar measure). The rate of convergence is governed by the representation theory of $G$; see Diaconis \cite{Diaconis}.

\subsection*{Summary: the exponential map across groups}

\begin{center}
\renewcommand{\arraystretch}{1.5}
\small
\begin{tabular}{lclll}
\textbf{Lie algebra $\mathfrak{g}$} & $\xrightarrow{\exp}$ & \textbf{Lie group $G$} & \textbf{$\exp$ global?} & \textbf{CLT limit} \\
\hline
$(\R, +)$ & $\to$ & $(\Rpos, \times)$ & Yes (iso.) & Gaussian $\to$ log-normal \\
$\mathfrak{so}(3) \cong \R^3$ & Rodrigues & $SO(3)$ & No (periodic) & short-time Gaussian \\
$\mathfrak{su}(2) \cong \R^3$ & Euler & $SU(2) \cong S^3$ & No (compact) & short-time Gaussian \\
$\mathfrak{sl}(n)$ & $\to$ & $SL(n)$ & No (local) & Gaussian on $\mathfrak{sl}(n)$
\end{tabular}
\end{center}

\noindent In every case, the exponential map $\exp : \mathfrak{g} \to G$ serves as the bridge between the additive (Lie algebra) and multiplicative (Lie group) formulations of the CLT. The abelian case $(\R, +) \to (\Rpos, \times)$ is the unique instance where this bridge is a global isomorphism, making the additive--multiplicative duality exact at all scales.

\subsection*{Spectral decomposition and irreversible convergence}

The preceding discussion treats the non-abelian extension at the level of the exponential map (algebra $\to$ group). We now develop the \emph{spectral} side: how the heat semigroup decomposes via representation theory, and how convergence rates connect back to the Wold decomposition and Szeg\H{o} theory of \cref{sec:hardy}.

\begin{theorem}[Peter-Weyl {\cite[Theorem~4.28]{Hall}}]\label{thm:peter-weyl}
Let $G$ be a compact Lie group, and let $\widehat{G}$ denote the set of equivalence classes of irreducible unitary representations $(\rho, V_\rho)$ of $G$. Write $d_\rho = \dim V_\rho$.
\begin{enumerate}[label=(\roman*)]
\item \textbf{(Orthogonality.)} The matrix coefficients $\sqrt{d_\rho}\, \rho_{ij}$, ranging over all $[\rho] \in \widehat{G}$ and $1 \leq i, j \leq d_\rho$, form an orthonormal basis of $L^2(G)$ with respect to Haar measure.

\item \textbf{(Isometric decomposition.)} The non-abelian Fourier transform
\[
\widehat{f}(\rho) = \int_G f(g)\, \rho(g)^*\, dg \in \operatorname{End}(V_\rho)
\]
defines an isometry (Plancherel identity):
\[
\|f\|_{L^2(G)}^2 = \sum_{[\rho] \in \widehat{G}} d_\rho\, \|\widehat{f}(\rho)\|_{\mathrm{HS}}^2,
\]
where $\|\cdot\|_{\mathrm{HS}}$ is the Hilbert--Schmidt norm.

\item \textbf{(Convolution becomes multiplication.)} For $f_1, f_2 \in L^1(G)$,
\[
\widehat{f_1 * f_2}(\rho) = \widehat{f_1}(\rho)\, \widehat{f_2}(\rho).
\]
\end{enumerate}
\end{theorem}

\begin{proof}[Proof sketch]
Part~(i) follows from Schur orthogonality for compact groups. Part~(ii) is the Peter-Weyl theorem proper. Part~(iii) is immediate from Fubini and the representation property $\rho(gh) = \rho(g)\rho(h)$.
\end{proof}

For abelian $G$, each $V_\rho$ is one-dimensional ($d_\rho = 1$) and the matrix coefficient $\rho_{11}$ is a character $\chi \in \widehat{G}$. The Peter-Weyl decomposition reduces to the classical Fourier series ($G = \bT$) or Fourier transform ($G = \R$), recovering the framework of Sections~2--5. The non-abelian case is distinguished by $d_\rho > 1$: the Fourier transform is \emph{operator-valued}, and convolution becomes \emph{matrix} multiplication rather than scalar multiplication.

\begin{definition}[Casimir operator]\label{def:casimir}
Let $G$ be a compact Lie group with Lie algebra $\mathfrak{g}$, equipped with a bi-invariant inner product (e.g., the negative of the Killing form when $G$ is semisimple). Let $\{X_1, \ldots, X_d\}$ be an orthonormal basis of $\mathfrak{g}$. The \emph{Casimir operator} (Laplace--Beltrami operator) on $G$ is
\[
\Omega = -\sum_{i=1}^d X_i^2,
\]
where each $X_i$ acts on $C^\infty(G)$ as a left-invariant differential operator. By Schur's lemma \cite[Proposition~10.7]{Hall}, $\Omega$ acts on each irreducible representation $V_\rho$ as scalar multiplication:
\[
\Omega\big|_{V_\rho} = \lambda_\rho \cdot \operatorname{Id}_{V_\rho}, \qquad \lambda_\rho \geq 0,
\]
with $\lambda_\rho = 0$ if and only if $\rho$ is the trivial representation.
\end{definition}

This is the non-abelian generalization of the Fourier symbol $Q(\chi)$ from \cref{def:laplacian} and \cref{ex:laplacian-symbol}. On $\bT$, $\lambda_n = 4\pi^2 n^2$; on $SO(3)$, $\lambda_\ell = \ell(\ell+1)$ for the spin-$\ell$ representation.

\begin{proposition}[Heat kernel on compact groups]\label{prop:heat-compact}
Let $G$ be a compact Lie group with Casimir operator $\Omega$. The heat kernel $p_t \in C^\infty(G)$, the fundamental solution of $\partial_t u = -\Omega u$, decomposes as
\[
p_t(g) = \sum_{[\rho] \in \widehat{G}} d_\rho\, e^{-\lambda_\rho t}\, \chi_\rho(g), \qquad t > 0,
\]
where $\chi_\rho(g) = \operatorname{tr}(\rho(g))$ is the character of~$\rho$. The series converges absolutely and uniformly for $t > 0$.
\end{proposition}

\begin{proof}[Proof sketch]
Expand $p_t$ via \cref{thm:peter-weyl}. Since $\Omega$ acts on each $V_\rho$ as $\lambda_\rho \cdot \operatorname{Id}$, the heat equation becomes $\partial_t \widehat{p_t}(\rho) = -\lambda_\rho\, \widehat{p_t}(\rho)$ with solution $\widehat{p_t}(\rho) = e^{-\lambda_\rho t}\, \operatorname{Id}_{V_\rho}$. Taking the trace and summing over $\widehat{G}$ gives the character expansion. Absolute convergence follows from the Weyl dimension formula $d_\rho = O(\lambda_\rho^{d/2})$, ensuring the series converges for $t > 0$; see \cite[\S 3]{Diaconis}.
\end{proof}

\begin{corollary}[Spectral gap and convergence rate]\label{cor:spectral-gap}
Let $\lambda_1 > 0$ be the spectral gap of $\Omega$ (the smallest nonzero Casimir eigenvalue). For any initial density $\rho_0 \in L^2(G)$ with $\int_G \rho_0 = 1$,
\[
\left\| \rho_t - \frac{1}{\operatorname{vol}(G)} \right\|_{L^2(G)} \leq e^{-\lambda_1 t}\left\| \rho_0 - \frac{1}{\operatorname{vol}(G)} \right\|_{L^2(G)}.
\]
Every square-integrable distribution on $G$ converges to Haar measure at exponential rate~$\lambda_1$. For the heat kernel ($\rho_0 = \delta_e$, not in $L^2$), the Plancherel series $\sum_{\rho \neq \mathbf{1}} d_\rho^2\, e^{-2\lambda_\rho t}$ converges for each $t > 0$ and decays at rate $e^{-2\lambda_1 t}$.
\end{corollary}

\begin{proof}
Write $f_0 = \rho_0 - 1/\operatorname{vol}(G)$. By Peter--Weyl, $f_0 = \sum_{\rho \neq \mathbf{1}} \widehat{f_0}(\rho)$ with $\|f_0\|_2^2 = \sum_{\rho \neq \mathbf{1}} d_\rho\, \|\widehat{f_0}(\rho)\|_{\mathrm{HS}}^2$. The heat semigroup gives $\widehat{f_t}(\rho) = e^{-\lambda_\rho t}\, \widehat{f_0}(\rho)$, so $\|f_t\|_2^2 = \sum_{\rho \neq \mathbf{1}} d_\rho\, \|\widehat{f_0}(\rho)\|_{\mathrm{HS}}^2\, e^{-2\lambda_\rho t} \leq e^{-2\lambda_1 t}\, \|f_0\|_2^2$.
\end{proof}

This is the non-abelian analogue of \cref{thm:heat-semigroup}: on $\bT$, $\lambda_1 = 4\pi^2$ and the convergence rate of the Jacobi theta function to~$1$ is $e^{-4\pi^2 t}$ (cf.\ \cref{ex:gaussian-instances}). On $(\R, +)$ (non-compact), the spectral gap vanishes and convergence is polynomial ($t^{-1/2}$)---precisely the CLT rate.

\Cref{prop:heat-compact} treats pure diffusion ($b = 0$). Adding a drift field yields the Fokker--Planck equation on~$G$, whose Peter--Weyl decomposition reveals a fundamental rigidity: diffusion dominates drift in every representation block.

\begin{proposition}[Fokker--Planck equation on compact groups]\label{prop:fpe-compact}
Let $G$ be a compact Lie group with Casimir operator $\Omega$, and let $b : G \to \mathfrak{g}$ be a smooth drift field. The density $\rho_t$ (with respect to Haar measure) of the stochastic process
\[
dg_t = g_t \cdot \bigl(b(g_t)\, dt + \sqrt{2\varepsilon}\, dW_t\bigr), \qquad \varepsilon > 0,
\]
satisfies the Fokker--Planck equation
\begin{equation}\label{eq:fpe-group}
\partial_t \rho = -\varepsilon\, \Omega\rho - \operatorname{div}_G(\rho\, b),
\end{equation}
where $\operatorname{div}_G$ is the Riemannian divergence on~$G$.
Under the Peter--Weyl decomposition (writing $\pi$ for irreducible
representations to free $\rho$ for the density):
\begin{enumerate}[label=(\roman*)]
\item \textbf{Left-invariant drift} ($b \in \mathfrak{g}$ constant). Since left-invariant vector fields are divergence-free with respect to Haar measure, \eqref{eq:fpe-group} becomes $\partial_t \rho = -\varepsilon\, \Omega\rho - X_b\rho$. Taking the non-abelian Fourier transform:
\[
\partial_t \widehat{\rho}(\pi) = \bigl(-\varepsilon\, \lambda_\pi\, \operatorname{Id}_{V_\pi} - d\pi(b)\bigr)\, \widehat{\rho}(\pi), \qquad [\pi] \in \widehat{G}.
\]
Each irreducible block evolves as an independent $d_\pi \times d_\pi$ matrix ODE. Since $\pi$ is unitary, $d\pi(b)$ is skew-Hermitian, so every eigenvalue of the coefficient matrix has real part $-\varepsilon\lambda_\pi$. The drift rotates the density within each block but cannot prevent diffusive decay.

\item \textbf{General drift} ($b(g)$ state-dependent). The product $\rho \cdot b$ mixes Peter--Weyl blocks via the Clebsch--Gordan coefficients: the evolution of $\widehat{\rho}(\pi)$ depends on $\widehat{\rho}(\sigma)$ for $\sigma \neq \pi$.

\item \textbf{Pure diffusion} ($b = 0$). Recovers \cref{prop:heat-compact}: $\widehat{\rho}_t(\pi) = e^{-\varepsilon\lambda_\pi t}\, \widehat{\rho}_0(\pi)$.
\end{enumerate}
In particular, for left-invariant drift, $\|\widehat{\rho}_t(\pi)\|_{\mathrm{HS}} \leq e^{-\varepsilon\lambda_\pi t}\, \|\widehat{\rho}_0(\pi)\|_{\mathrm{HS}}$ for each~$\pi$, so the spectral gap bound of \cref{cor:spectral-gap} is robust: convergence to Haar measure at rate $\varepsilon\lambda_1$ persists regardless of the drift.
\end{proposition}

\begin{proof}[Proof sketch]
The generator of the diffusion is $\mathcal{L}f = -\varepsilon\,\Omega f + X_b f$; the FPE is its $L^2(G)$-adjoint. For left-invariant~$b$, the Fourier transform of $X_b\rho$ is computed by integration by parts against the matrix coefficients: using $X_b[\pi(g)^*] = -d\pi(b)\,\pi(g)^*$ (which follows from unitarity of~$\pi$), one obtains $\widehat{X_b\rho}(\pi) = d\pi(b)\,\widehat{\rho}(\pi)$. The skew-Hermiticity of $d\pi(b)$ then gives the decay estimate via $\operatorname{Re}\,\sigma(-\varepsilon\lambda_\pi\,\operatorname{Id} - d\pi(b)) = \{-\varepsilon\lambda_\pi\}$.
\end{proof}

\begin{remark}[Irreversibility: the non-abelian Wold analogue]
The heat semigroup $\{e^{-t\Omega}\}_{t \geq 0}$ on $L^2(G)$ is a contraction semigroup: each representation mode decays monotonically as $e^{-\lambda_\rho t}$, and $\lambda_\rho > 0$ for all non-trivial~$\rho$. This gives the non-abelian analogue of the information loss in the Wold decomposition (\cref{thm:wold}):
\[
\renewcommand{\arraystretch}{1.4}
\begin{array}{lll}
\hline
& \textbf{Abelian (\S\ref{sec:hardy})} & \textbf{Non-abelian} \\
\hline
\text{State space} & L^2(\bT) & L^2(G) \\
\text{Decomposition} & \Htwo \oplus (\Htwo)^\perp & \bigoplus_\rho d_\rho\, \operatorname{End}(V_\rho) \\
\text{Forward operator} & \text{Shift } S & \text{Heat semigroup } e^{-t\Omega} \\
\text{Information loss} & \ker(S^*) = \operatorname{span}\{1\} & \text{Decay rate } \lambda_\rho \text{ per mode} \\
\text{Deterministic part} & \mathcal{H}_u = \{0\} \text{ (Szeg\H{o})} & \text{Trivial rep.\ } \rho = \mathbf{1} \text{ (Haar)} \\
\text{CLT fixed point} & \text{Gaussian density} & \text{Haar measure} \\
\hline
\end{array}
\]
In both cases, the ``deterministic'' component---the part surviving forward evolution---is the fixed point of the respective semigroup. The Szeg\H{o} condition $\int \log f > -\infty$ (ensuring $\mathcal{H}_u = \{0\}$ in the abelian Wold theorem) is replaced by the spectral gap condition $\lambda_1 > 0$ (ensuring all non-trivial modes decay).
\end{remark}

\begin{remark}[Non-compact groups: Plancherel theorem and Mackey's machine]
For non-compact $G$, the Peter-Weyl theorem is replaced by the Plancherel theorem: the sum over $\widehat{G}$ becomes an integral with respect to Plancherel measure $\mu_{\mathrm{Pl}}$ on the tempered dual $\widehat{G}_{\mathrm{temp}}$ \cite[Chapter~7]{Folland}:
\[
\|f\|_{L^2(G)}^2 = \int_{\widehat{G}_{\mathrm{temp}}} \|\widehat{f}(\rho)\|_{\mathrm{HS}}^2\, d\mu_{\mathrm{Pl}}(\rho).
\]
For semidirect products $G = K \ltimes V$ ($K$ compact, $V$ a vector group), Mackey's theory of induced representations \cite[Chapter~6]{Folland} parametrizes $\widehat{G}$ by orbits of $K$ acting on $\widehat{V}$:
\begin{enumerate}[label=(\roman*)]
\item Each orbit $\mathcal{O} \subset \widehat{V}$ with stabilizer $K_\xi$ (the little group at $\xi \in \mathcal{O}$) and irreducible $\sigma$ of $K_\xi$ yields an irreducible $\operatorname{Ind}_{K_\xi \ltimes V}^G(\sigma \otimes e^{i\xi})$ of~$G$.
\item The Casimir spectrum has both discrete (from $K$) and continuous (from $V$) components.
\item The spectral gap may vanish: $\Omega$ on $L^2(G)$ typically has continuous spectrum down to~$0$, so convergence to equilibrium is polynomial ($t^{-d/2}$) rather than exponential.
\end{enumerate}
\textbf{Example: $\boldsymbol{SE(3) = SO(3) \ltimes \R^3}$.} Here $K = SO(3)$ acts on $\widehat{V} = \R^3$ by rotation. The orbits are spheres $S^2_r$ ($r > 0$) and the origin. The little group at $\xi = (0,0,r)$ is $SO(2)$. Each irreducible of $SE(3)$ is induced from $SO(2) \ltimes \R^3$, parametrized by $(r, m)$ with $r > 0$, $m \in \Z$. The Casimir spectrum is continuous in~$r$ and discrete in~$m$: no spectral gap, polynomial convergence.
\end{remark}

\begin{remark}[Representation-theoretic decomposition and optimal control]
The Peter-Weyl decomposition has a precise but limited connection to optimal control on Lie groups. Consider the controlled diffusion on a compact~$G$:
\[
dg_t = g_t \cdot \bigl(u(t)\, dt + \sqrt{2\varepsilon}\, dW_t\bigr), \qquad u(t) \in \mathfrak{g},
\]
with cost $J = \mathbb{E}\bigl[\int_0^T L(g_t, u_t)\, dt + \Phi(g_T)\bigr]$. The Hamilton--Jacobi--Bellman equation for the value function $V \in C^\infty(G)$ is
\[
\partial_t V - \varepsilon\, \Omega V + \min_{u \in \mathfrak{g}}\bigl\{L(g, u) + \langle dV,\, g \cdot u\rangle\bigr\} = 0.
\]
In the Peter-Weyl basis, the three terms decompose differently:
\begin{enumerate}[label=(\roman*)]
\item The diffusion term $-\varepsilon\, \Omega V$ \textbf{block-diagonalizes}: $\Omega$ acts as $\lambda_\rho$ on each $V_\rho$ (\cref{def:casimir}).
\item A left-invariant drift $g \cdot X$ acts on $V_\rho$ as the derived representation $d\rho(X)$, a $d_\rho \times d_\rho$ matrix. The free (uncontrolled) dynamics within each block is a \textbf{linear matrix ODE}.
\item The optimization $\min_u$ and a general cost $L(g, u)$ create \textbf{inter-block coupling}: if $L$ is not a class function (constant on conjugacy classes), its Peter-Weyl expansion mixes different representations.
\end{enumerate}
For the special case of left-invariant dynamics with quadratic cost on $\mathfrak{g}$ (the Lie-group analogue of LQR, linear-quadratic regulation), the HJB equation decouples into independent finite-dimensional Riccati equations per representation block. For general nonlinear costs, Peter-Weyl provides an optimal spectral Galerkin basis respecting the group symmetry, but the $\min_u$ coupling remains. The spectral gap $\lambda_1$ (\cref{cor:spectral-gap}) provides an analytic lower bound on the convergence rate of policy evaluation, independently of the policy.
\end{remark}

\begin{remark}[Rigid-body configuration spaces and spectral truncation]
The configuration space of a floating-base rigid body with $n$ revolute joints is $Q = SE(3) \times \bT^n$. The FPE--HJB framework (\cref{prop:fpe-compact} and the preceding remark) extends to such products: the Peter--Weyl basis for $L^2(G_1 \times G_2)$ consists of tensor products $\rho_1 \otimes \rho_2$ with Casimir eigenvalue $\lambda_{\rho_1 \otimes \rho_2} = \lambda_{\rho_1} + \lambda_{\rho_2}$.

Two obstructions arise in high dimension. First, $SE(3)$ is non-compact: its Plancherel decomposition has continuous spectrum and no spectral gap, as detailed in the preceding remark on $SE(3)$. Second, on the compact factor $SO(3) \times \bT^n$, the Weyl eigenvalue asymptotics give $O(\Lambda^{(3+n)/2})$ spectral degrees of freedom below a cutoff~$\Lambda$---a curse of dimensionality for any spectral Galerkin truncation of the FPE or HJB equation. Whether the kinematic chain structure (nearest-neighbor coupling in the joint topology, which enforces sparsity in the Peter--Weyl expansion of physical Hamiltonians) can mitigate this growth remains open.
Alternatively, stochastic methods bypass the spectral curse entirely. Path sampling of the SDE in \cref{prop:fpe-compact} yields Monte Carlo estimates of $\mathbb{E}[\phi(g_t)]$ at rate $O(N^{-1/2})$ for $N$ sample paths and bounded observables $\phi$, independently of $\dim G$---the CLT rate of the present paper. The spectral gap $\lambda_1$ (\cref{cor:spectral-gap}) re-enters as the mixing rate of the sampler, governing the number of effectively independent samples per unit time. In this sense, diffusion provides a \emph{universal} spectral regularization: modes with $\lambda_\rho \gg \varepsilon^{-1}t^{-1}$ are exponentially suppressed by the heat semigroup without being explicitly computed, and the operative complexity bound becomes $N^{-1/2}$ (stochastic, dimension-free) rather than $\Lambda^{(3+n)/2}$ (deterministic, dimension-cursed).
\end{remark}

\begin{remark}[Hopf--Cole linearization and path-integral control]
Specialize the HJB equation of the preceding remark to quadratic control cost
$L(g,u) = q(g) + \tfrac{1}{2}\|u\|_{\mathfrak{g}}^2$.
The $\min_u$ is a Legendre--Fenchel transform: the minimizer is
$u^* = -\nabla_{\mathfrak{g}} V$ with minimum value
$q(g) - \tfrac{1}{2}\|\nabla_{\mathfrak{g}} V\|^2$,
where $\nabla_{\mathfrak{g}} V$ denotes the gradient of $V$ with respect to
the inner product on $\mathfrak{g}$.
The specialized HJB becomes
\[
\partial_t V - \varepsilon\,\Omega V + q(g)
  - \tfrac{1}{2}\|\nabla_{\mathfrak{g}} V\|^2 = 0,
\quad V(T,g) = \Phi(g).
\]
In the deterministic limit $\varepsilon = 0$, the characteristics of the
resulting Hamilton--Jacobi equation are Euler--Poincar\'e trajectories on
$\mathfrak{g}$ (equivalently, Lie--Poisson flow on $\mathfrak{g}^*$ via the
inertia tensor $\mathbb{I}: \mathfrak{g} \to \mathfrak{g}^*$).

\medskip\noindent\textbf{Linearization.}
The SDE has noise covariance $\Sigma = 2\varepsilon\,\mathrm{Id}_{\mathfrak{g}}$
and the control cost has $R = \mathrm{Id}_{\mathfrak{g}}$.
Setting $\lambda = 2\varepsilon$ so that $\lambda R^{-1} = \Sigma$,
the Hopf--Cole substitution $\psi = e^{-V/(2\varepsilon)}$ exactly cancels the
quadratic gradient nonlinearity.
The derivation (chain rule for $\Omega(\log\psi)$, then
$\|\nabla_{\mathfrak{g}} V\|^2$ cancellation) yields the linearized PDE:
\[
\partial_t \psi + (-\varepsilon\,\Omega)\psi
  - \frac{q}{2\varepsilon}\,\psi = 0,
\quad \psi(T,g) = e^{-\Phi(g)/(2\varepsilon)}.
\]
This is a backward linear PDE---the heat equation on $G$ with killing
potential $q/(2\varepsilon)$.

\medskip\noindent\textbf{Feynman--Kac representation.}
Since $-\varepsilon\,\Omega$ is the generator of the uncontrolled diffusion
$dg_s = g_s \cdot \sqrt{2\varepsilon}\,dW_s$
(cf.\ proof sketch of \cref{prop:fpe-compact}):
\[
\psi(t,g) = \mathbb{E}_g\!\left[
  \exp\!\left(-\frac{1}{2\varepsilon}\int_t^T q(g_s)\,ds\right)
  e^{-\Phi(g_T)/(2\varepsilon)}
\right],
\]
where the expectation is over the uncontrolled (pure-diffusion) process
started at $g_t = g$.

\medskip\noindent\textbf{Optimal control recovery and MPPI.}
The optimal control is
$u^*(t,g) = 2\varepsilon\,\nabla_{\mathfrak{g}}\log\psi(t,g)$.
The nonlinear HJB is replaced by forward sampling of the SDE plus
exponential reweighting---this is model-predictive path-integral control
(MPPI) \cite{Kappen,Todorov}.
The Monte Carlo estimate
$\hat\psi = \frac{1}{N}\sum_{i=1}^N w^{(i)}$ with weights
$w^{(i)} = \exp\bigl(-\frac{1}{2\varepsilon}\int q\,ds
  - \frac{\Phi}{2\varepsilon}\bigr)$
converges at $O(N^{-1/2})$ for bounded costs, independently of
$\dim G$---the CLT rate of the present paper.

\medskip\noindent\textbf{Deterministic limit.}
As $\varepsilon \to 0$, a Schilder-type large-deviation argument (Laplace
approximation of the Feynman--Kac integral) shows
$V(t,g) \to \min_u\bigl\{\int_t^T L(g_s,u_s)\,ds + \Phi(g_T)\bigr\}$.
The minimizing path satisfies the Euler--Lagrange equation; on a Lie group
with left-invariant kinetic energy, this reduces to the Euler--Poincar\'e
equation on $\mathfrak{g}$.
\end{remark}

\begin{proposition}[Forward-only reduction via the exponential bridge]
\label{prop:forward-only}
Consider the MPPI estimator for the controlled diffusion on a Lie group~$G$
described in the preceding remark.
\begin{enumerate}[label=\textup{(\roman*)}]
\item \textbf{Arithmetic completeness.}
Every operation in the estimator is a finite composition of three primitives:
\begin{itemize}[nosep]
\item $\exp_G : \mathfrak{g} \to G$ \textup{(}Lie group exponential\textup{)},
\item $\log_G : G \to \mathfrak{g}$ \textup{(}Lie group logarithm\textup{)},
\item linear algebra \textup{(}matrix multiply-accumulate on
      $\mathfrak{g} \cong \R^d$\textup{)}.
\end{itemize}
No backward differentiation through the dynamics and no other transcendental
operations are required.
\item \textbf{Dimension-free convergence.}
The estimator converges at $O(N^{-1/2})$, independently of $\dim G$.
\end{enumerate}
\end{proposition}

\begin{proof}[Proof \textup{(}constructive enumeration\textup{)}]
We enumerate the four computational stages and verify that each uses only the
three primitives.

\medskip\noindent\emph{SDE step.}
$g_{t+h} = g_t \cdot \exp_G\!\bigl(h\,u + \sqrt{2\varepsilon h}\;\xi\bigr)$,
where $\xi \in \mathfrak{g}$ is Gaussian noise.
This requires one~$\exp_G$ and one matrix multiply-accumulate (MMA) in
$\mathfrak{g}$.

\medskip\noindent\emph{Cost evaluation.}
For a tracking cost $q(g) = \bigl\|\log_G(g^{-1}g_{\mathrm{target}})\bigr\|^2$:
one~$\log_G$ and one MMA (squared norm in~$\mathfrak{g}$).

\medskip\noindent\emph{Exponential reweighting.}
$w^{(i)} = \exp\!\bigl(-\mathrm{cost}^{(i)}/(2\varepsilon)\bigr)$:
one scalar exponential (the $\dim G = 1$ case of $\exp_G$).

\medskip\noindent\emph{Weighted average.}
$u^* = 2\varepsilon\,\sum_i w^{(i)}\xi^{(i)}\big/\sum_i w^{(i)}$:
MMA in~$\mathfrak{g}$.

\medskip\noindent
The weights $w^{(i)}$ are bounded (bounded cost, fixed~$\varepsilon$), so the
weighted average is a bounded-weight Monte Carlo estimator.  The CLT
(\cref{thm:add-clt,thm:mult-clt}) gives convergence at $O(N^{-1/2})$,
independently of~$\dim G$.
\end{proof}

\begin{remark}[The exponential bridge]
The proposition's three primitives $\{\exp,\log,\mathrm{MMA}\}$ are the
computational realization of a single algebraic fact---the exponential
homomorphism $\exp\colon (+) \to (\times)$---which appears in three guises
throughout this paper:
\begin{enumerate}[label=\textup{(\arabic*)}]
\item \emph{Fourier--Mellin duality}
      (\cref{thm:mellin-fourier}):
      $\log\colon (\Rpos, \times) \to (\R, +)$ converts multiplicative
      convolution to additive convolution.
\item \emph{Hopf--Cole linearization}
      (preceding remark):
      $\psi = e^{-V/(2\varepsilon)}$ converts the nonlinear HJB
      (quadratic $\|\nabla_{\mathfrak{g}} V\|^2$---multiplicative)
      to the linear heat equation
      (additive semigroup $e^{t\varepsilon\Omega}$).
\item \emph{Lie group exponential}:
      $\exp_G\colon \mathfrak{g} \to G$ converts the linear/additive
      Lie algebra to the nonlinear/multiplicative group.
      The Lie--Euler integrator traverses this bridge at every time step.
\end{enumerate}
\end{remark}

\begin{remark}[Generalized path integral and Riccati compression]
The acronym GPI in \cref{thm:in-context-gpi} admits a second reading:
\emph{Generalized Path Integral}.
The two readings coincide.
In the filtration $\mathbb{F}\colon (\mathbb{N},\leq) \to \Proj_Y$
of \cref{thm:in-context-gpi}, each inclusion
$\mathcal{G}_t \hookrightarrow \mathcal{G}_{t+1}$ extends the
$\sigma$-algebra by one datum---a batch of observations in $\R^d$.
The conditional expectation $E[Y \mid \mathcal{G}_t]$ is the current
estimate; the residual $R_Y(\mathcal{G}_t) = \|Y - E[Y \mid
\mathcal{G}_t]\|_2^2$ is the unexplained variance.
This is the evaluation half.

The MPPI estimator of \cref{prop:forward-only} supplies the
improvement half: sample $N$ trajectories, assign Boltzmann weights
$w^{(i)} \propto \exp\!\bigl(-J^{(i)}/\varepsilon\bigr)$, and compute
the weighted average.
The Boltzmann reweighting is exactly the exponential bridge applied to
the path measure---it converts the uniform (prior) measure over
trajectories to the optimal (posterior) measure, with the cost
functional $J$ playing the role of the negative log-likelihood.
Each reweighted sample extends the filtration; convergence at
$O(N^{-1/2})$ is guaranteed by the CLT
(\cref{thm:add-clt,thm:mult-clt}), independently of $\dim G$.

\medskip\noindent\textbf{Riccati compression.}
Specialize to $G = (\R^n, +)$ with linear dynamics
$\dot x = Ax + Bu$ and quadratic cost
$\int_0^T (x^\top Q\,x + u^\top R\,u)\,dt$
(the LQG regime---linear dynamics, quadratic cost, Gaussian noise),
where $A \in \R^{n \times n}$, $B \in \R^{n \times m}$,
$Q \succeq 0$, and $R \succ 0$.
The Feynman--Kac path integral is Gaussian and evaluates in closed
form; the Boltzmann-weighted average compresses to the matrix Riccati
ODE:
\[
\dot P = A^\top\! P + PA - PBR^{-1}B^\top\! P + Q,
\]
where $P(t) \in \R^{n \times n}$ is the cost-to-go matrix and
$V(x) = x^\top P\, x$ the value function.
No sampling is needed: the exact solution replaces the Monte Carlo
estimator entirely, and the filtration convergence of
\cref{thm:in-context-gpi} is achieved analytically.
The dual equation (Kalman--Bucy filter) compresses the estimation
filtration in the same way.
The CLT rate $O(N^{-1/2})$ of \cref{prop:forward-only} thus quantifies
the cost of leaving the Riccati regime: on a nonlinear Lie group~$G$
with non-quadratic cost, the path integral cannot be evaluated
analytically and one pays $O(N^{-1/2})$ per sample for the Monte Carlo
approximation.
\end{remark}

\subsection*{Articulated rigid-body dynamics and the Lie-algebraic factorization}

An articulated mechanism---a tree of $n+1$ rigid bodies connected by joints,
with body~$0$ as the floating base---provides a concrete test of the framework
above.
The configuration space is $Q = SE(3) \times \bT^n$, exactly the product
considered in the preceding remark, and every step of the classical recursive
dynamics algorithms is an operation on $\mathfrak{se}(3)$ or its dual
\cite{Featherstone,MurrayLiSastry}.

\medskip\noindent\textbf{(A) Configuration space and product of exponentials.}
Label the bodies $0, 1, \ldots, n$ along the kinematic tree.
Body $i$ has joint angle $\theta_i \in \bT$ and screw axis
$S_i \in \mathfrak{se}(3)$ (in body frame).
The pose of body $i$ relative to the world frame is given by the product of
exponentials:
\[
T_i(q)
= g_0 \, \bar{M}_1 \, e^{S_1 \theta_1} \cdots \bar{M}_i \, e^{S_i \theta_i},
\]
where $g_0 \in SE(3)$ is the base pose and
$\bar{M}_i \in SE(3)$ is the reference-configuration offset of joint $i$.

\medskip\noindent\textbf{(B) Recursive velocity propagation---the adjoint
representation in action.}
The body velocity of link $i$ is
$\mathcal{V}_i^b = T_i^{-1}\dot{T}_i \in \mathfrak{se}(3)$, and it
satisfies the recursion
\[
\mathcal{V}_i^b
= \mathrm{Ad}_{(e^{S_i\theta_i}\bar{M}_i)^{-1}}\,\mathcal{V}_{i-1}^b
  + S_i\,\dot{\theta}_i.
\]
The recursion is built entirely from the \textbf{adjoint representation}
$\mathrm{Ad}_g : \mathfrak{se}(3) \to \mathfrak{se}(3)$---the same object
governing the Peter--Weyl decomposition of $L^2(SE(3))$.

\medskip\noindent\textbf{(C) Spatial inertia and the Euler--Poincar\'e
equation.}
Each body $i$ carries a spatial inertia
$\cM_i \in \mathrm{Sym}^+(\mathfrak{se}(3)^*)$: a $6 \times 6$
positive-definite matrix in Pl\"ucker coordinates encoding mass, center of
mass, and rotational inertia.
The Newton--Euler equation for body $i$ in body frame is
\[
\mathcal{F}_i
= \cM_i\,\dot{\mathcal{V}}_i
  + \mathrm{ad}^*_{\mathcal{V}_i}\!(\cM_i\,\mathcal{V}_i),
\]
where $\mathrm{ad}^*$ is the coadjoint action of $\mathfrak{se}(3)$ on
$\mathfrak{se}(3)^*$.
This is exactly the \textbf{Euler--Poincar\'e equation} on $\mathfrak{se}(3)$
for a single rigid body.
For a free rigid body ($\mathcal{F} = 0$), writing
$\mu = \cM\,\mathcal{V} \in \mathfrak{se}(3)^*$ gives
$\dot{\mu} = \mathrm{ad}^*_{\mathcal{V}}\mu$; the rotational component
recovers \textbf{Euler's equation} $\dot{\Pi} = \Pi \times \omega$.

\medskip\noindent\textbf{(D) Recursive Newton--Euler Algorithm
(RNEA)---$O(n)$ inverse dynamics.}
\begin{itemize}[nosep]
\item \emph{Forward pass} (root $\to$ leaves): propagate velocities
  $\mathcal{V}_i$ and accelerations $\dot{\mathcal{V}}_i$ using $\mathrm{Ad}$.
\item \emph{Backward pass} (leaves $\to$ root): compute body wrenches
  $\mathcal{F}_i$ via Euler--Poincar\'e, propagate to parents using the
  coadjoint $\mathrm{Ad}^*$, and project onto joint axes:
  $\tau_i = S_i^T \mathcal{F}_i$.
\end{itemize}
Complexity: $O(n)$---one Euler--Poincar\'e evaluation and one $\mathrm{Ad}$
transform per body.

\noindent In Lagrangian coordinates \cite{CarpentierMansard}, the inverse dynamics reads
\begin{equation}\label{eq:lagrangian-eom}
\tau = M(q)\,\ddot{q} + C(q,\dot{q})\,\dot{q} + g(q)
  - \sum_i J_i(q)^T f_i^{\mathrm{ext}},
\end{equation}
where $M(q) \in \R^{n \times n}$ is the joint-space inertia matrix,
$C(q,\dot{q})\,\dot{q}$ collects Coriolis and centrifugal terms,
$g(q)$ is the gravity vector, and $J_i(q)$ is the frame Jacobian at
which external force $f_i^{\mathrm{ext}}$ is applied.
RNEA evaluates the right-hand side in $O(n)$; its analytical
derivatives $\partial\tau/\partial q$, $\partial\tau/\partial\dot{q}$
are likewise $O(n)$ \cite{CarpentierMansard}.

\medskip\noindent\textbf{(E) Articulated Body Algorithm (ABA)---$O(n)$
forward dynamics.}
The ABA computes the forward dynamics---inverting \eqref{eq:lagrangian-eom}
to obtain $\ddot{q}$ given $(\tau, f^{\mathrm{ext}})$---without forming
$M$ explicitly.
The key object is the \emph{articulated-body inertia} $\cM_i^A$, the
effective inertia of body $i$ and all its descendants accounting for joint
freedom.
The recursion (leaves $\to$ root) is
\[
\cM_i^A = \cM_i + \!\!\sum_{j \in \mathrm{children}(i)}
  \mathrm{Ad}_{T_{ji}}^*
  \!\left(\cM_j^A
    - \cM_j^A S_j \bigl(S_j^T \cM_j^A S_j\bigr)^{\!-1} S_j^T \cM_j^A
  \right) \mathrm{Ad}_{T_{ji}}.
\]
The inner expression
$\cM^A - \cM^A S(S^T\cM^A S)^{-1} S^T \cM^A$ is a \textbf{Schur
complement}---the projection of the inertia onto the space orthogonal to the
joint axis.
This is a Riccati recursion on $\mathrm{Sym}^+(\mathfrak{se}(3)^*)$.
Complexity: $O(n)$.

\medskip\noindent\textbf{(E$'$) Contact forces: recover then eliminate.}
When the mechanism interacts with the environment, unilateral contacts
impose constraints $\phi_k(q) \geq 0$ with associated Jacobians
$J_k^c = \partial\phi_k/\partial q$.
The contact forces $\lambda_k \geq 0$ satisfy the complementarity
condition $0 \leq \lambda_k \perp \phi_k(q) \geq 0$ \cite{MenagerCarpentier2025}.
In the Lagrangian form \eqref{eq:lagrangian-eom}, they enter as
additional external forces $f_k^{\mathrm{ext}} = \lambda_k\,\hat{n}_k$
at the contact points.

\medskip\noindent\emph{Recover.}
Substituting the forward dynamics into the contact constraint
$\ddot{\phi}_k \geq 0$ yields the \textbf{Delassus system}:
\[
G\,\lambda = J^c\, M(q)^{-1}\bigl(\tau - C\dot{q} - g\bigr)
  + \dot{J}^c\,\dot{q},
\qquad G = J^c\, M(q)^{-1} (J^c)^T.
\]
The Delassus operator $G$ is a Schur complement of the mass matrix
onto the contact space---the same algebraic structure as the
articulated-body inertia projection in part~(E).

\medskip\noindent\emph{Eliminate.}
Once $\lambda$ is recovered, substitute back:
$\ddot{q} = M^{-1}\bigl(\tau - C\dot{q} - g + (J^c)^T\lambda\bigr)$.
The contact forces are eliminated from the forward dynamics; what
remains is branchless arithmetic on $\mathfrak{se}(3)$.

\medskip\noindent\emph{Barrier smoothing.}
The complementarity condition $0 \leq \lambda \perp \phi \geq 0$ is a
hard branch.
Replacing it with a log-barrier potential
$V_{\mathrm{barrier}}(\phi) = -\varepsilon\log\phi$ yields a smooth
contact force $\lambda_\varepsilon = \varepsilon / \phi$ that
approximates the complementarity as $\varepsilon \to 0$.
(For exact treatment without relaxation, the complementarity can be
handled as a quadratic program with complementarity
constraints~\cite{MenagerCarpentier2025}.)
The ``recover'' step becomes differentiable, and the MPPI estimator
of \cref{prop:forward-only} samples through the barrier-smoothed
landscape without branch.

\medskip\noindent\textbf{(F) Lie--Poisson duality: velocity
$\leftrightarrow$ momentum.}
The Legendre transform $\mathbb{FL}: \mathfrak{g} \to \mathfrak{g}^*$ via the
composite inertia tensor maps body velocities to spatial momenta
\cite{MarsdenRatiu}.
Euler--Poincar\'e on $\bigoplus_i \mathfrak{se}(3)$ (velocity variables) is
dual to Lie--Poisson on $\bigoplus_i \mathfrak{se}(3)^*$ (momentum
variables):
RNEA computes the Euler--Poincar\'e side; ABA computes the Lie--Poisson side.
The inertia tensor
$\cM_i : \mathfrak{se}(3) \to \mathfrak{se}(3)^*$ is the bridge, exactly as
in the classical duality $\mathfrak{g} \leftrightarrow \mathfrak{g}^*$ of
geometric mechanics.

\medskip\noindent\textbf{(G) Connection to the paper's spectral and
stochastic framework.}
Three connections tie the recursive dynamics to the results of this paper:

\begin{enumerate}[label=(\roman*)]
\item \emph{Representation-theoretic.}
The $\mathrm{Ad}$ and $\mathrm{Ad}^*$ in Featherstone's recursions are the
adjoint and coadjoint representations of $SE(3)$.
In the Peter--Weyl/Plancherel decomposition of $L^2(SE(3))$, these same
representations govern how Fourier modes transform under group action.
The tree topology constrains them to act only between parent--child
pairs---the ``nearest-neighbor coupling'' that enforces sparsity in the
Peter--Weyl expansion, as noted in the preceding remark.

\item \emph{Riccati structure.}
The ABA recursion is a discrete matrix Riccati equation on
$\mathfrak{se}(3)^*$.
The per-block Riccati of the HJB remark (for left-invariant quadratic cost)
is the spectral-domain counterpart: ABA factors the Riccati in the spatial
domain (tree topology), while Peter--Weyl factors it in the spectral domain
(representation blocks).

\item \emph{Stochastic extension and MPPI.}
Adding noise to joint torques or base forces gives an SDE on
$SE(3) \times \bT^n$.
The FPE has product Casimir (\cref{prop:fpe-compact}).
The Hopf--Cole linearization (preceding remark) converts the HJB for the
articulated system to a Feynman--Kac expectation.
Each forward sample of the SDE can be computed in $O(n)$ per time step using
the recursive structure---the tree factorization is preserved under
randomization.
The Monte Carlo estimator converges at the CLT rate $O(N^{-1/2})$,
independently of $n$ and $\dim G = 6 + n$:
Featherstone provides the $O(n)$ forward simulator; the CLT provides the
dimension-free convergence guarantee.
This is the MPPI framework for articulated rigid-body control.
\end{enumerate}

\subsection*{Free probability}
In free probability, the role of the Gaussian is taken by the \emph{semicircle distribution} (Wigner), and the free multiplicative analogue is the \emph{Marchenko--Pastur distribution}. The R-transform and S-transform serve as the free analogues of the Fourier and Mellin transforms, respectively, linearizing free additive and free multiplicative convolution. The duality between additive and multiplicative free convolutions mirrors the classical Fourier--Mellin duality studied here, but in a non-commutative setting; see Voiculescu--Dykema--Nica \cite{VDN}.

%% ============================================================
\begin{thebibliography}{99}

\bibitem{Aubin}
J.-P.~Aubin,
\emph{Viability Theory}, new ed.,
Modern Birkh\"auser Classics, Birkh\"auser, 2009.

\bibitem{BQ}
Y.~Benoist and J.-F.~Quint,
\emph{Random Walks on Reductive Groups},
Ergebnisse der Mathematik und ihrer Grenzgebiete, vol.~62, Springer, 2016.

\bibitem{Billingsley}
P.~Billingsley,
\emph{Probability and Measure}, 3rd ed.,
Wiley, New York, 1995.

\bibitem{BottcherSilbermann}
A.~B\"ottcher and B.~Silbermann,
\emph{Analysis of Toeplitz Operators}, 2nd ed.,
Springer Monographs in Mathematics, Springer, 2006.

\bibitem{CarpentierMansard}
J.~Carpentier and N.~Mansard.
\newblock Analytical derivatives of rigid body dynamics algorithms.
\newblock In \emph{Robotics: Science and Systems~(RSS)}, 2018.

\bibitem{MenagerCarpentier2025}
E.~M\'enager, A.~Bambade, W.~Jallet, A.~De~Marchi, and J.~Carpentier.
\newblock Contact-implicit inverse dynamics.
\newblock Submitted to \emph{IEEE Trans.\ Robotics}, 2025.

\bibitem{Diaconis}
P.~Diaconis,
Group representations in probability and statistics,
\emph{IMS Lecture Notes---Monograph Series}, vol.~11, 1988.

\bibitem{Douglas}
R.~G.~Douglas,
\emph{Banach Algebra Techniques in Operator Theory}, 2nd ed.,
Graduate Texts in Mathematics, vol.~179, Springer, 1998.

\bibitem{Durrett}
R.~Durrett,
\emph{Probability: Theory and Examples}, 5th ed.,
Cambridge University Press, 2019.

\bibitem{Featherstone}
R.~Featherstone,
\emph{Rigid Body Dynamics Algorithms},
Springer, 2008.

\bibitem{Folland}
G.~B.~Folland,
\emph{A Course in Abstract Harmonic Analysis}, 2nd ed.,
CRC Press, 2015.

\bibitem{Garnett}
J.~B.~Garnett,
\emph{Bounded Analytic Functions}, revised 1st ed.,
Graduate Texts in Mathematics, vol.~236, Springer, 2007.

\bibitem{Hall}
B.~C.~Hall,
\emph{Lie Groups, Lie Algebras, and Representations: An Elementary Introduction}, 2nd ed.,
Graduate Texts in Mathematics, vol.~222, Springer, 2015.

\bibitem{Hoffman}
K.~Hoffman,
\emph{Banach Spaces of Analytic Functions},
Prentice-Hall, 1962; reprinted Dover, 1988.

\bibitem{Heyer}
H.~Heyer,
\emph{Probability Measures on Locally Compact Groups},
Ergebnisse der Mathematik und ihrer Grenzgebiete, vol.~94, Springer, 1977.

\bibitem{HewittRoss}
E.~Hewitt and K.~A.~Ross,
\emph{Abstract Harmonic Analysis}, vol.~I,
Grundlehren der mathematischen Wissenschaften, vol.~115, Springer, 1963.

\bibitem{Kappen}
H.~J.~Kappen,
Path integrals and symmetry breaking for optimal control theory,
\emph{J.\ Stat.\ Mech.}, 2005:P11011, 2005.

\bibitem{MarsdenRatiu}
J.~E.~Marsden and T.~S.~Ratiu,
\emph{Introduction to Mechanics and Symmetry}, 2nd ed.,
Texts in Applied Mathematics, vol.~17, Springer, 1999.

\bibitem{MochizukiIII}
S.~Mochizuki,
Inter-universal Teichm\"uller theory~III: Canonical splittings of the log-theta-lattice,
\emph{Publ.\ Res.\ Inst.\ Math.\ Sci.} \textbf{57} (2021), 723--1185.

\bibitem{MurrayLiSastry}
R.~M.~Murray, Z.~Li, and S.~S.~Sastry,
\emph{A Mathematical Introduction to Robotic Manipulation},
CRC Press, 1994.

\bibitem{Nikolski}
N.~K.~Nikolski,
\emph{Operators, Functions, and Systems: An Easy Reading}, vols.~1--2,
Mathematical Surveys and Monographs, AMS, 2002.

\bibitem{Parthasarathy}
K.~R.~Parthasarathy,
\emph{Probability Measures on Metric Spaces},
Academic Press, 1967; reprinted AMS Chelsea, 2005.

\bibitem{Rudin}
W.~Rudin,
\emph{Fourier Analysis on Groups},
Wiley-Interscience, New York, 1962.

\bibitem{ScholzeStix}
P.~Scholze and J.~Stix,
Why abc is still a conjecture,
preprint, 2018.

\bibitem{Simon}
B.~Simon,
\emph{Orthogonal Polynomials on the Unit Circle}, 2 vols.,
AMS Colloquium Publications, vol.~54, 2005.

\bibitem{SzNF}
B.~Sz.-Nagy, C.~Foias, H.~Bercovici, and L.~K\'erchy,
\emph{Harmonic Analysis of Operators on Hilbert Space}, 2nd ed.,
Universitext, Springer, 2010.

\bibitem{Todorov}
E.~Todorov,
Linearly-solvable Markov decision problems,
in \emph{Advances in Neural Information Processing Systems~19} (NIPS),
2006, pp.~1369--1376.

\bibitem{VDN}
D.~V.~Voiculescu, K.~J.~Dykema, and A.~Nica,
\emph{Free Random Variables},
CRM Monograph Series, vol.~1, AMS, 1992.

\end{thebibliography}

\end{document}
